% Assessment of computational methods for diabetic retinopathy lesion segmentation in fundus images
% Layout follows the Nature Methods "Analysis" format of the template paper
% (Lu et al., Nat. Methods 23, 248-259): Results, Discussion, Methods, Extended Data.
%
% Build:  make            (figures, tables, numbers, bibliography, PDF)
%         make final      (same, with every number printed as measured)
\documentclass[10pt,a4paper]{article}

\usepackage[T1]{fontenc}
\usepackage[utf8]{inputenc}
\usepackage{mathptmx}
\usepackage[scaled=0.9]{helvet}
\usepackage{courier}
\usepackage[a4paper,margin=20mm,top=22mm,bottom=24mm]{geometry}
\usepackage{microtype}
\usepackage{amsmath,amssymb}
\usepackage{graphicx}
\usepackage[table]{xcolor}
\usepackage{booktabs,colortbl,multirow,makecell,array}
\usepackage{threeparttable}
\usepackage{adjustbox}
\usepackage{pifont}
\usepackage[font={small,sf},labelfont=bf,labelsep=space,justification=justified,singlelinecheck=false]{caption}
\usepackage{titlesec}
\usepackage{enumitem}
\usepackage[most]{tcolorbox}
\usepackage[switch]{lineno}
\usepackage{placeins}
\usepackage[numbers,super,sort&compress]{natbib}
\usepackage[colorlinks=true,allcolors=linkblue]{hyperref}

\definecolor{linkblue}{HTML}{1C5CAB}
% Shared table macros. \input this after loading:
%   booktabs, colortbl, xcolor (with the "table" option), multirow, makecell,
%   threeparttable, adjustbox, pifont, amssymb, array.
%
% Draft and final. In draft, every number that is still simulated is typeset
% in \phcolor so it cannot be mistaken for a result. \retifinaltrue prints
% all numbers in black.
\newif\ifretifinal
\retifinalfalse

\definecolor{phcolor}{HTML}{5B7FA6}
\definecolor{sdgray}{HTML}{8A8A86}
\definecolor{rulegray}{HTML}{C9C8C0}
\definecolor{headgray}{HTML}{F4F3EF}
\definecolor{inkmuted}{HTML}{6B6B66}
\definecolor{famgeneral}{HTML}{2A78D6}
\definecolor{famauto}{HTML}{1BAF7A}
\definecolor{famfundus}{HTML}{EB6834}
\definecolor{famfoundation}{HTML}{4A3AA7}
\definecolor{tierfine}{HTML}{1D3F5E}
\definecolor{tiermixed}{HTML}{3F6F9C}
\definecolor{tiercoarse}{HTML}{7AA3C7}
\definecolor{tierregion}{HTML}{AABFD3}
\definecolor{deltaup}{HTML}{0B7A3E}
\definecolor{deltadown}{HTML}{B3261E}
\definecolor{barfill}{HTML}{8FB0CC}

% --- numbers
\newcommand{\ph}[1]{\ifretifinal#1\else{\color{phcolor}#1}\fi}          % simulated value
\newcommand{\best}[1]{\textbf{#1}}                                      % best in column
\newcommand{\second}[1]{\underline{#1}}                                 % second best
\newcommand{\sd}[1]{{\scriptsize\color{sdgray}\,$\pm$\,#1}}             % seed s.d.
\newcommand{\up}[1]{{\scriptsize\color{deltaup}$\blacktriangle$\,#1}}   % gain
\newcommand{\down}[1]{{\scriptsize\color{deltadown}$\blacktriangledown$\,#1}} % loss
\newcommand{\same}[1]{{\scriptsize\color{sdgray}#1}}
\newcommand{\na}{{\color{rulegray}--}}
\newcommand{\tbd}{{\color{deltadown}\scriptsize tbd}}

% --- marks
\newcommand{\yes}{{\color{deltaup}\ding{51}}}
\newcommand{\no}{{\color{rulegray}\ding{55}}}
\newcommand{\famchip}[1]{{\color{fam#1}\rule[0.05ex]{0.62em}{0.62em}}}       % family chip
\newcommand{\tier}[1]{{\color{tier#1}\rule[0.05ex]{0.62em}{0.62em}}}     % granularity chip
\newcommand{\tierbadge}[2]{{\setlength{\fboxsep}{1.6pt}\colorbox{tier#1}{\color{white}\scriptsize\sffamily\strut #2}}}
\newcommand{\rlevel}[1]{% reproducibility meter, #1 = 0..4
  \ifcase#1 $\circ\circ\circ\circ$\or $\bullet\circ\circ\circ$\or $\bullet\bullet\circ\circ$\or
  $\bullet\bullet\bullet\circ$\or $\bullet\bullet\bullet\bullet$\fi}
\newcommand{\inlinebar}[1]{{\color{barfill}\rule[0.1ex]{#1em}{0.55em}}}  % #1 = length in em

% --- structure
\newcommand{\grouprow}[2]{\rowcolor{headgray}\multicolumn{#1}{l}{\textbf{#2}}\\}
\newcommand{\colh}[1]{\textbf{#1}}
\newcommand{\tablefont}{\footnotesize\setlength{\tabcolsep}{4.2pt}\renewcommand{\arraystretch}{1.12}}
\newcommand{\phnote}{\ifretifinal\else\item[] {\color{phcolor}Values in this colour are simulated placeholders, not results.}\fi}

\ifdefined\retisegfinal\retifinaltrue\fi   % "make final" sets this; or write \retifinaltrue here

% --- numbers quoted in the text: \rn{key}, defined in generated/numbers.tex
\makeatletter
\newcommand{\rn}[1]{\@ifundefined{rn@#1}{\textbf{??#1??}}{\@nameuse{rn@#1}}}
\makeatother
% Generated by scripts/make_numbers.py. Do not edit by hand.
\makeatletter
\@namedef{rn@backbone.best}{HRNet-W48}
\@namedef{rn@backbone.best.ddr}{HRNet-W48}
\@namedef{rn@backbone.span}{0.05--0.15}
\@namedef{rn@bstd.DDR.EX.best}{0.65}
\@namedef{rn@bstd.DDR.EX.median}{0.60}
\@namedef{rn@bstd.DDR.HE.best}{0.56}
\@namedef{rn@bstd.DDR.HE.median}{0.46}
\@namedef{rn@bstd.DDR.MA.best}{0.19}
\@namedef{rn@bstd.DDR.MA.median}{0.15}
\@namedef{rn@bstd.DDR.SE.best}{0.53}
\@namedef{rn@bstd.DDR.SE.median}{0.39}
\@namedef{rn@bstd.DDR.best.model}{SAMed}
\@namedef{rn@bstd.DDR.best.value}{0.46}
\@namedef{rn@bstd.DDR.gap12}{0.010}
\@namedef{rn@bstd.DDR.sd.max}{0.027}
\@namedef{rn@bstd.DDR.sd.median}{0.009}
\@namedef{rn@bstd.DDR.second.model}{RETFound + UPerNet}
\@namedef{rn@bstd.DDR.second.value}{0.45}
\@namedef{rn@bstd.DDR.unet.SE}{0.38}
\@namedef{rn@bstd.DDR.worst.model}{FCT}
\@namedef{rn@bstd.DDR.worst.value}{0.18}
\@namedef{rn@bstd.IDRiD.EX.best}{0.89}
\@namedef{rn@bstd.IDRiD.EX.median}{0.87}
\@namedef{rn@bstd.IDRiD.HE.best}{0.74}
\@namedef{rn@bstd.IDRiD.HE.median}{0.66}
\@namedef{rn@bstd.IDRiD.MA.best}{0.50}
\@namedef{rn@bstd.IDRiD.MA.median}{0.44}
\@namedef{rn@bstd.IDRiD.SE.best}{0.78}
\@namedef{rn@bstd.IDRiD.SE.median}{0.68}
\@namedef{rn@bstd.IDRiD.best.model}{SAMed}
\@namedef{rn@bstd.IDRiD.best.value}{0.71}
\@namedef{rn@bstd.IDRiD.gap12}{0.000}
\@namedef{rn@bstd.IDRiD.sd.max}{0.084}
\@namedef{rn@bstd.IDRiD.sd.median}{0.010}
\@namedef{rn@bstd.IDRiD.second.model}{SAM-Adapter}
\@namedef{rn@bstd.IDRiD.second.value}{0.71}
\@namedef{rn@bstd.IDRiD.unet.SE}{0.68}
\@namedef{rn@bstd.IDRiD.worst.model}{FCT}
\@namedef{rn@bstd.IDRiD.worst.value}{0.30}
\@namedef{rn@cd.chi2}{97.9}
\@namedef{rn@cd.first.model}{SAMed}
\@namedef{rn@cd.first.rank}{1.63}
\@namedef{rn@cd.last.model}{FCT}
\@namedef{rn@cd.last.rank}{17.88}
\@namedef{rn@cd.ntop}{11}
\@namedef{rn@cd.p}{$<$ 0.001}
\@namedef{rn@cd.second.model}{HRNet-W48}
\@namedef{rn@cd.second.rank}{4.13}
\@namedef{rn@cd.third.model}{Swin-Unet}
\@namedef{rn@cd.third.rank}{5.25}
\@namedef{rn@cd.value}{9.90}
\@namedef{rn@ci.DDR.EX.width}{0.14}
\@namedef{rn@ci.DDR.HE.width}{0.19}
\@namedef{rn@ci.DDR.MA.width}{0.07}
\@namedef{rn@ci.DDR.SE.width}{0.24}
\@namedef{rn@ci.DDR.width.hi}{0.19}
\@namedef{rn@ci.DDR.width.lo}{0.07}
\@namedef{rn@ci.IDRiD.EX.width}{0.12}
\@namedef{rn@ci.IDRiD.HE.width}{0.11}
\@namedef{rn@ci.IDRiD.MA.width}{0.08}
\@namedef{rn@ci.IDRiD.SE.width}{0.33}
\@namedef{rn@ci.IDRiD.width.hi}{0.12}
\@namedef{rn@ci.IDRiD.width.lo}{0.08}
\@namedef{rn@ci.hrnet.DDR.HE}{0.477 (0.39--0.59)}
\@namedef{rn@ci.hrnet.DDR.MA}{0.155 (0.12--0.20)}
\@namedef{rn@ci.hrnet.DDR.SE}{0.550 (0.42--0.65)}
\@namedef{rn@ci.hrnet.IDRiD.HE}{0.684 (0.63--0.72)}
\@namedef{rn@ci.hrnet.IDRiD.MA}{0.446 (0.41--0.48)}
\@namedef{rn@ci.hrnet.IDRiD.SE}{0.783 (0.59--0.87)}
\@namedef{rn@ci.nnunet.DDR.HE}{0.501 (0.39--0.57)}
\@namedef{rn@ci.nnunet.DDR.MA}{0.154 (0.11--0.17)}
\@namedef{rn@ci.nnunet.DDR.SE}{0.458 (0.30--0.56)}
\@namedef{rn@ci.nnunet.IDRiD.HE}{0.671 (0.61--0.72)}
\@namedef{rn@ci.nnunet.IDRiD.MA}{0.499 (0.42--0.51)}
\@namedef{rn@ci.nnunet.IDRiD.SE}{0.641 (0.40--0.79)}
\@namedef{rn@ci.retfound.DDR.HE}{0.491 (0.39--0.60)}
\@namedef{rn@ci.retfound.DDR.MA}{0.189 (0.14--0.24)}
\@namedef{rn@ci.retfound.DDR.SE}{0.546 (0.38--0.66)}
\@namedef{rn@ci.retfound.IDRiD.HE}{0.718 (0.66--0.75)}
\@namedef{rn@ci.retfound.IDRiD.MA}{0.484 (0.45--0.52)}
\@namedef{rn@ci.retfound.IDRiD.SE}{0.734 (0.50--0.85)}
\@namedef{rn@ci.samed.DDR.HE}{0.574 (0.52--0.64)}
\@namedef{rn@ci.samed.DDR.MA}{0.173 (0.13--0.22)}
\@namedef{rn@ci.samed.DDR.SE}{0.460 (0.30--0.60)}
\@namedef{rn@ci.samed.IDRiD.HE}{0.736 (0.67--0.78)}
\@namedef{rn@ci.samed.IDRiD.MA}{0.449 (0.40--0.50)}
\@namedef{rn@ci.samed.IDRiD.SE}{0.790 (0.61--0.88)}
\@namedef{rn@clin.ddr.hi.auc}{0.75}
\@namedef{rn@clin.ddr.hi.auc.n}{12}
\@namedef{rn@clin.ddr.hi.maupr}{0.36}
\@namedef{rn@clin.ddr.hi.n}{203}
\@namedef{rn@clin.ddr.lo.auc}{0.54}
\@namedef{rn@clin.ddr.lo.auc.n}{12}
\@namedef{rn@clin.ddr.lo.maupr}{0.21}
\@namedef{rn@clin.ddr.lo.n}{20}
\@namedef{rn@clin.fgadr.hi.auc}{0.76}
\@namedef{rn@clin.fgadr.hi.auc.n}{19}
\@namedef{rn@clin.fgadr.hi.maupr}{0.22}
\@namedef{rn@clin.fgadr.hi.n}{459}
\@namedef{rn@clin.fgadr.lo.auc}{0.71}
\@namedef{rn@clin.fgadr.lo.auc.n}{19}
\@namedef{rn@clin.fgadr.lo.maupr}{0.04}
\@namedef{rn@clin.fgadr.lo.n}{94}
\@namedef{rn@clin.maples.hi.auc}{0.83}
\@namedef{rn@clin.maples.hi.auc.n}{19}
\@namedef{rn@clin.maples.hi.maupr}{0.40}
\@namedef{rn@clin.maples.hi.n}{4}
\@namedef{rn@clin.maples.lo.auc}{0.73}
\@namedef{rn@clin.maples.lo.auc.n}{19}
\@namedef{rn@clin.maples.lo.maupr}{0.33}
\@namedef{rn@clin.maples.lo.n}{56}
\@namedef{rn@clin.retles.hi.auc}{0.83}
\@namedef{rn@clin.retles.hi.auc.n}{19}
\@namedef{rn@clin.retles.hi.maupr}{0.17}
\@namedef{rn@clin.retles.hi.n}{328}
\@namedef{rn@clin.retles.lo.auc}{0.81}
\@namedef{rn@clin.retles.lo.auc.n}{19}
\@namedef{rn@clin.retles.lo.maupr}{0.07}
\@namedef{rn@clin.retles.lo.n}{150}
\@namedef{rn@cost.ddr.max}{72}
\@namedef{rn@cost.ddr.max.model}{FCT}
\@namedef{rn@cost.ddr.ratio}{3.9}
\@namedef{rn@cost.max}{20}
\@namedef{rn@cost.max.model}{FCT}
\@namedef{rn@cost.min}{1.5}
\@namedef{rn@cost.min.model}{HRDecoder}
\@namedef{rn@cross.fgadr.fine}{0.15}
\@namedef{rn@cross.max}{0.28}
\@namedef{rn@cross.mean}{0.16}
\@namedef{rn@e1.m2mrf.DDR.exact}{0.401}
\@namedef{rn@e1.m2mrf.DDR.fn}{0.492}
\@namedef{rn@e1.m2mrf.DDR.gap}{0.091}
\@namedef{rn@e1.m2mrf.DDR.paper}{0.489}
\@namedef{rn@e1.m2mrf.DDR.stored}{0.426}
\@namedef{rn@e1.m2mrf.IDRiD.exact}{0.639}
\@namedef{rn@e1.m2mrf.IDRiD.fn}{0.676}
\@namedef{rn@e1.m2mrf.IDRiD.gap}{0.036}
\@namedef{rn@e1.m2mrf.IDRiD.paper}{0.672}
\@namedef{rn@e1.m2mrf.IDRiD.stored}{0.647}
\@namedef{rn@e1r.HRDecoder.IDRiD.delta}{$-$0.041}
\@namedef{rn@e1r.HRDecoder.IDRiD.exact}{0.674}
\@namedef{rn@e1r.HRDecoder.IDRiD.fn}{0.672}
\@namedef{rn@e1r.M2MRF.IDRiD.delta}{$-$0.060}
\@namedef{rn@e1r.M2MRF.IDRiD.exact}{0.549}
\@namedef{rn@e1r.M2MRF.IDRiD.fn}{0.612}
\@namedef{rn@e1r.first.exact}{0.493}
\@namedef{rn@e1r.first.fn}{0.576}
\@namedef{rn@e4.fgadr}{+0.02}
\@namedef{rn@e4.fgadr383}{+0.01}
\@namedef{rn@e4.fgadronly}{$-$0.16}
\@namedef{rn@e4.ref}{0.56}
\@namedef{rn@e4.retles}{$-$0.04}
\@namedef{rn@e4.retlesonly}{$-$0.35}
\@namedef{rn@e7.aupr.1440}{0.42}
\@namedef{rn@e7.aupr.512}{0.38}
\@namedef{rn@e7.aupr.native}{0.38}
\@namedef{rn@e7.small.1440}{0.28}
\@namedef{rn@e7.small.512}{0.15}
\@namedef{rn@e8.ddr.10}{0.31}
\@namedef{rn@e8.ddr.100}{0.38}
\@namedef{rn@e8.ddr.50}{0.36}
\@namedef{rn@e8.ddr.all}{0.37}
\@namedef{rn@e8.tjdr.100}{0.50}
\@namedef{rn@e8.tjdr.fine}{0.56}
\@namedef{rn@e9.range}{0.14}
\@namedef{rn@family.auto.independent}{0.36}
\@namedef{rn@family.auto.internal}{0.54}
\@namedef{rn@family.foundation.independent}{0.38}
\@namedef{rn@family.foundation.internal}{0.55}
\@namedef{rn@family.fundus.independent}{0.32}
\@namedef{rn@family.fundus.internal}{0.47}
\@namedef{rn@family.general.independent}{0.30}
\@namedef{rn@family.general.internal}{0.47}
\@namedef{rn@fm.any.ref.nnunet.IDRiD}{0.78}
\@namedef{rn@fm.any.sam.DDR}{0.03}
\@namedef{rn@fm.any.sam.IDRiD}{0.07}
\@namedef{rn@fm.any.samtwo.DDR}{0.03}
\@namedef{rn@fm.any.samtwo.IDRiD}{0.08}
\@namedef{rn@fm.box.medsam.DDR}{0.65}
\@namedef{rn@fm.box.medsam.IDRiD}{0.66}
\@namedef{rn@fm.box.sam.DDR}{0.77}
\@namedef{rn@fm.box.sam.IDRiD}{0.86}
\@namedef{rn@fm.box.samtwo.DDR}{0.76}
\@namedef{rn@fm.box.samtwo.IDRiD}{0.85}
\@namedef{rn@fm.boxref.m2mrf.DDR}{0.67}
\@namedef{rn@fm.boxref.m2mrf.IDRiD}{0.81}
\@namedef{rn@fm.boxref.nnunet.DDR}{0.71}
\@namedef{rn@fm.boxref.nnunet.IDRiD}{0.86}
\@namedef{rn@fm.enc.c1.ind}{-0.026}
\@namedef{rn@fm.enc.c1.int}{-0.053}
\@namedef{rn@fm.enc.c1.int.abs}{0.053}
\@namedef{rn@fm.enc.c1.std.abs}{0.045}
\@namedef{rn@fm.enc.c2.ind}{0.006}
\@namedef{rn@fm.enc.c2.int}{-0.020}
\@namedef{rn@fm.enc.c2.int.abs}{0.020}
\@namedef{rn@fm.enc.c2.std.abs}{0.003}
\@namedef{rn@fm.enc.di.TJDR-std}{0.62}
\@namedef{rn@fm.enc.di.TJDR-uwf}{0.14}
\@namedef{rn@fm.enc.di.ind}{0.41}
\@namedef{rn@fm.enc.di.int}{0.51}
\@namedef{rn@fm.enc.dinov2.ind}{0.41}
\@namedef{rn@fm.enc.i.c1.ind}{-0.038}
\@namedef{rn@fm.enc.i.c1.int}{-0.021}
\@namedef{rn@fm.enc.i.c1.int.abs}{0.021}
\@namedef{rn@fm.enc.i.c2.ind}{-0.006}
\@namedef{rn@fm.enc.i.c2.int}{-0.028}
\@namedef{rn@fm.enc.i.c2.int.abs}{0.028}
\@namedef{rn@fm.enc.i.di.MAPLES-DR}{0.44}
\@namedef{rn@fm.enc.i.di.TJDR-std}{0.53}
\@namedef{rn@fm.enc.i.di.TJDR-uwf}{0.25}
\@namedef{rn@fm.enc.i.di.ind}{0.41}
\@namedef{rn@fm.enc.i.di.int}{0.74}
\@namedef{rn@fm.enc.i.mi.MAPLES-DR}{0.43}
\@namedef{rn@fm.enc.i.mi.TJDR-std}{0.54}
\@namedef{rn@fm.enc.i.mi.TJDR-uwf}{0.27}
\@namedef{rn@fm.enc.i.mi.ind}{0.41}
\@namedef{rn@fm.enc.i.mi.int}{0.73}
\@namedef{rn@fm.enc.i.rd.MAPLES-DR}{0.42}
\@namedef{rn@fm.enc.i.rd.TJDR-std}{0.56}
\@namedef{rn@fm.enc.i.rd.TJDR-uwf}{0.22}
\@namedef{rn@fm.enc.i.rd.ind}{0.40}
\@namedef{rn@fm.enc.i.rd.int}{0.71}
\@namedef{rn@fm.enc.i.rm.MAPLES-DR}{0.41}
\@namedef{rn@fm.enc.i.rm.TJDR-std}{0.46}
\@namedef{rn@fm.enc.i.rm.TJDR-uwf}{0.25}
\@namedef{rn@fm.enc.i.rm.ind}{0.37}
\@namedef{rn@fm.enc.i.rm.int}{0.71}
\@namedef{rn@fm.enc.imagenet.ind}{0.41}
\@namedef{rn@fm.enc.imagenet.int}{0.52}
\@namedef{rn@fm.enc.mi.TJDR-std}{0.63}
\@namedef{rn@fm.enc.mi.TJDR-uwf}{0.13}
\@namedef{rn@fm.enc.mi.ind}{0.41}
\@namedef{rn@fm.enc.mi.int}{0.52}
\@namedef{rn@fm.enc.rd.TJDR-std}{0.62}
\@namedef{rn@fm.enc.rd.TJDR-uwf}{0.15}
\@namedef{rn@fm.enc.rd.ind}{0.42}
\@namedef{rn@fm.enc.rd.int}{0.49}
\@namedef{rn@fm.enc.retfound.ind}{0.42}
\@namedef{rn@fm.enc.retfound.int}{0.49}
\@namedef{rn@fm.enc.rm.TJDR-std}{0.59}
\@namedef{rn@fm.enc.rm.TJDR-uwf}{0.11}
\@namedef{rn@fm.enc.rm.ind}{0.39}
\@namedef{rn@fm.enc.rm.int}{0.47}
\@namedef{rn@fm.enc.sets}{the standard-field TJDR set, the ultra-widefield TJDR set and MAPLES-DR}
\@namedef{rn@fm.oracle.all}{0.85}
\@namedef{rn@fm.oracle.five}{0.35}
\@namedef{rn@fm.oracle.one}{0.16}
\@namedef{rn@fm.oracle.twenty}{0.55}
\@namedef{rn@fm.peft.model}{SAMed}
\@namedef{rn@fm.peft.value}{0.46}
\@namedef{rn@fm.point.all.DDR}{0.02}
\@namedef{rn@fm.point.all.IDRiD}{0.12}
\@namedef{rn@fm.ref.nnunet}{0.43}
\@namedef{rn@fm.sam3.EX.IDRiD}{0.41}
\@namedef{rn@fm.sam3.EX.max}{0.41}
\@namedef{rn@fm.sam3.EX.min}{0.03}
\@namedef{rn@fm.sam3.HE.IDRiD}{0.29}
\@namedef{rn@fm.sam3.HE.max}{0.29}
\@namedef{rn@fm.sam3.HE.min}{0.02}
\@namedef{rn@fm.sam3.MA.IDRiD}{0.01}
\@namedef{rn@fm.sam3.MA.max}{0.01}
\@namedef{rn@fm.sam3.MA.min}{0.00}
\@namedef{rn@fm.sam3.SE.IDRiD}{0.21}
\@namedef{rn@fm.sam3.SE.max}{0.21}
\@namedef{rn@fm.sam3.SE.min}{0.00}
\@namedef{rn@fm.wording.ratio}{27.9}
\@namedef{rn@fm.zero.model}{SAM 3}
\@namedef{rn@fm.zero.ub.model}{SAM 2}
\@namedef{rn@fm.zero.ub.value}{0.03}
\@namedef{rn@fm.zero.value}{0.40}
\@namedef{rn@funnel.code}{6}
\@namedef{rn@funnel.dr}{21}
\@namedef{rn@funnel.full}{2}
\@namedef{rn@funnel.train}{4}
\@namedef{rn@human.expert.MA.hi}{0.022}
\@namedef{rn@human.expert.MA.lo}{0.001}
\@namedef{rn@human.expert.auc.hi}{0.93}
\@namedef{rn@human.expert.auc.lo}{0.87}
\@namedef{rn@human.expert.maupr.hi}{0.28}
\@namedef{rn@human.expert.maupr.lo}{0.15}
\@namedef{rn@human.expert.maupr.median}{0.21}
\@namedef{rn@human.models.DDR.auc.median}{0.84}
\@namedef{rn@human.models.DDR.auc.nreach}{7}
\@namedef{rn@human.models.DDR.maupr.median}{0.28}
\@namedef{rn@human.models.IDRiD.auc.median}{0.83}
\@namedef{rn@human.models.IDRiD.auc.nreach}{5}
\@namedef{rn@human.models.IDRiD.maupr.median}{0.25}
\@namedef{rn@human.models.MA.median}{0.039}
\@namedef{rn@human.models.auc.hi}{0.94}
\@namedef{rn@human.models.auc.lo}{0.71}
\@namedef{rn@human.models.below}{3}
\@namedef{rn@human.models.maupr.hi}{0.32}
\@namedef{rn@human.models.maupr.lo}{0.09}
\@namedef{rn@human.models.pairs}{38}
\@namedef{rn@impl.max}{0.060}
\@namedef{rn@impl.mean}{0.005}
\@namedef{rn@impl.moved}{10}
\@namedef{rn@impl.tau}{0.88}
\@namedef{rn@indep.DDR.DiaRetDB1.first.intrank}{10}
\@namedef{rn@indep.DDR.DiaRetDB1.first.model}{DeepLabv3}
\@namedef{rn@indep.DDR.DiaRetDB1.first.value}{0.348}
\@namedef{rn@indep.DDR.DiaRetDB1.mean}{0.251}
\@namedef{rn@indep.DDR.DiaRetDB1.pct}{$-$32}
\@namedef{rn@indep.DDR.DiaRetDB1.rho}{0.78}
\@namedef{rn@indep.DDR.FGADR.drop}{0.16}
\@namedef{rn@indep.DDR.FGADR.first.intrank}{1}
\@namedef{rn@indep.DDR.FGADR.first.model}{SAMed}
\@namedef{rn@indep.DDR.FGADR.first.value}{0.288}
\@namedef{rn@indep.DDR.FGADR.gain}{-0.16}
\@namedef{rn@indep.DDR.FGADR.mean}{0.211}
\@namedef{rn@indep.DDR.FGADR.nlower}{19}
\@namedef{rn@indep.DDR.FGADR.p}{$<$ 0.001}
\@namedef{rn@indep.DDR.FGADR.pct}{$-$43}
\@namedef{rn@indep.DDR.FGADR.rho}{0.79}
\@namedef{rn@indep.DDR.IDRiD.first.intrank}{3}
\@namedef{rn@indep.DDR.IDRiD.first.model}{HRNet-W48}
\@namedef{rn@indep.DDR.IDRiD.first.value}{0.643}
\@namedef{rn@indep.DDR.IDRiD.mean}{0.550}
\@namedef{rn@indep.DDR.IDRiD.pct}{+49}
\@namedef{rn@indep.DDR.IDRiD.rho}{0.92}
\@namedef{rn@indep.DDR.MAPLES-DR.drop}{-0.00}
\@namedef{rn@indep.DDR.MAPLES-DR.first.intrank}{1}
\@namedef{rn@indep.DDR.MAPLES-DR.first.model}{SAMed}
\@namedef{rn@indep.DDR.MAPLES-DR.first.value}{0.470}
\@namedef{rn@indep.DDR.MAPLES-DR.gain}{0.00}
\@namedef{rn@indep.DDR.MAPLES-DR.mean}{0.372}
\@namedef{rn@indep.DDR.MAPLES-DR.nlower}{6}
\@namedef{rn@indep.DDR.MAPLES-DR.p}{= 0.595}
\@namedef{rn@indep.DDR.MAPLES-DR.pct}{+1}
\@namedef{rn@indep.DDR.MAPLES-DR.rho}{0.82}
\@namedef{rn@indep.DDR.Retinal-Lesions.drop}{0.18}
\@namedef{rn@indep.DDR.Retinal-Lesions.first.intrank}{10}
\@namedef{rn@indep.DDR.Retinal-Lesions.first.model}{DeepLabv3}
\@namedef{rn@indep.DDR.Retinal-Lesions.first.value}{0.282}
\@namedef{rn@indep.DDR.Retinal-Lesions.gain}{-0.18}
\@namedef{rn@indep.DDR.Retinal-Lesions.mean}{0.193}
\@namedef{rn@indep.DDR.Retinal-Lesions.nlower}{19}
\@namedef{rn@indep.DDR.Retinal-Lesions.p}{$<$ 0.001}
\@namedef{rn@indep.DDR.Retinal-Lesions.pct}{$-$48}
\@namedef{rn@indep.DDR.Retinal-Lesions.rho}{0.76}
\@namedef{rn@indep.DDR.TJDR-std.drop}{-0.12}
\@namedef{rn@indep.DDR.TJDR-std.first.intrank}{3}
\@namedef{rn@indep.DDR.TJDR-std.first.model}{HRNet-W48}
\@namedef{rn@indep.DDR.TJDR-std.first.value}{0.616}
\@namedef{rn@indep.DDR.TJDR-std.gain}{0.12}
\@namedef{rn@indep.DDR.TJDR-std.mean}{0.490}
\@namedef{rn@indep.DDR.TJDR-std.nlower}{0}
\@namedef{rn@indep.DDR.TJDR-std.p}{$<$ 0.001}
\@namedef{rn@indep.DDR.TJDR-std.pct}{+33}
\@namedef{rn@indep.DDR.TJDR-std.rho}{0.88}
\@namedef{rn@indep.DDR.TJDR-uwf.drop}{0.20}
\@namedef{rn@indep.DDR.TJDR-uwf.first.intrank}{9}
\@namedef{rn@indep.DDR.TJDR-uwf.first.model}{MedSAM (full FT)}
\@namedef{rn@indep.DDR.TJDR-uwf.first.value}{0.243}
\@namedef{rn@indep.DDR.TJDR-uwf.gain}{-0.20}
\@namedef{rn@indep.DDR.TJDR-uwf.mean}{0.173}
\@namedef{rn@indep.DDR.TJDR-uwf.nlower}{19}
\@namedef{rn@indep.DDR.TJDR-uwf.p}{$<$ 0.001}
\@namedef{rn@indep.DDR.TJDR-uwf.pct}{$-$53}
\@namedef{rn@indep.DDR.TJDR-uwf.rho}{0.01}
\@namedef{rn@indep.DDR.fine.first.model}{SAMed}
\@namedef{rn@indep.DDR.fine.first.value}{0.426}
\@namedef{rn@indep.DDR.fine.hrnet}{0.412}
\@namedef{rn@indep.DDR.int.mean}{0.369}
\@namedef{rn@indep.DDR.ndrop}{14}
\@namedef{rn@indep.DDR.rho}{0.80}
\@namedef{rn@indep.DDR.sametop}{the same model}
\@namedef{rn@indep.DDR.top.model}{SAMed}
\@namedef{rn@indep.DDR.top.value}{0.43}
\@namedef{rn@indep.IDRiD.DDR.first.intrank}{1}
\@namedef{rn@indep.IDRiD.DDR.first.model}{SAMed}
\@namedef{rn@indep.IDRiD.DDR.first.value}{0.368}
\@namedef{rn@indep.IDRiD.DDR.mean}{0.286}
\@namedef{rn@indep.IDRiD.DDR.pct}{$-$55}
\@namedef{rn@indep.IDRiD.DDR.rho}{0.75}
\@namedef{rn@indep.IDRiD.DiaRetDB1.first.intrank}{18}
\@namedef{rn@indep.IDRiD.DiaRetDB1.first.model}{DeepLabv3}
\@namedef{rn@indep.IDRiD.DiaRetDB1.first.value}{0.311}
\@namedef{rn@indep.IDRiD.DiaRetDB1.mean}{0.239}
\@namedef{rn@indep.IDRiD.DiaRetDB1.pct}{$-$62}
\@namedef{rn@indep.IDRiD.DiaRetDB1.rho}{0.48}
\@namedef{rn@indep.IDRiD.FGADR.drop}{0.43}
\@namedef{rn@indep.IDRiD.FGADR.first.intrank}{1}
\@namedef{rn@indep.IDRiD.FGADR.first.model}{SAMed}
\@namedef{rn@indep.IDRiD.FGADR.first.value}{0.280}
\@namedef{rn@indep.IDRiD.FGADR.gain}{-0.43}
\@namedef{rn@indep.IDRiD.FGADR.mean}{0.196}
\@namedef{rn@indep.IDRiD.FGADR.nlower}{19}
\@namedef{rn@indep.IDRiD.FGADR.p}{$<$ 0.001}
\@namedef{rn@indep.IDRiD.FGADR.pct}{$-$69}
\@namedef{rn@indep.IDRiD.FGADR.rho}{0.53}
\@namedef{rn@indep.IDRiD.MAPLES-DR.drop}{0.28}
\@namedef{rn@indep.IDRiD.MAPLES-DR.first.intrank}{10}
\@namedef{rn@indep.IDRiD.MAPLES-DR.first.model}{TC-Net}
\@namedef{rn@indep.IDRiD.MAPLES-DR.first.value}{0.412}
\@namedef{rn@indep.IDRiD.MAPLES-DR.gain}{-0.28}
\@namedef{rn@indep.IDRiD.MAPLES-DR.mean}{0.350}
\@namedef{rn@indep.IDRiD.MAPLES-DR.nlower}{19}
\@namedef{rn@indep.IDRiD.MAPLES-DR.p}{$<$ 0.001}
\@namedef{rn@indep.IDRiD.MAPLES-DR.pct}{$-$44}
\@namedef{rn@indep.IDRiD.MAPLES-DR.rho}{0.71}
\@namedef{rn@indep.IDRiD.Retinal-Lesions.drop}{0.46}
\@namedef{rn@indep.IDRiD.Retinal-Lesions.first.intrank}{18}
\@namedef{rn@indep.IDRiD.Retinal-Lesions.first.model}{DeepLabv3}
\@namedef{rn@indep.IDRiD.Retinal-Lesions.first.value}{0.260}
\@namedef{rn@indep.IDRiD.Retinal-Lesions.gain}{-0.46}
\@namedef{rn@indep.IDRiD.Retinal-Lesions.mean}{0.172}
\@namedef{rn@indep.IDRiD.Retinal-Lesions.nlower}{19}
\@namedef{rn@indep.IDRiD.Retinal-Lesions.p}{$<$ 0.001}
\@namedef{rn@indep.IDRiD.Retinal-Lesions.pct}{$-$73}
\@namedef{rn@indep.IDRiD.Retinal-Lesions.rho}{0.52}
\@namedef{rn@indep.IDRiD.TJDR-std.drop}{0.21}
\@namedef{rn@indep.IDRiD.TJDR-std.first.intrank}{1}
\@namedef{rn@indep.IDRiD.TJDR-std.first.model}{SAMed}
\@namedef{rn@indep.IDRiD.TJDR-std.first.value}{0.534}
\@namedef{rn@indep.IDRiD.TJDR-std.gain}{-0.21}
\@namedef{rn@indep.IDRiD.TJDR-std.mean}{0.424}
\@namedef{rn@indep.IDRiD.TJDR-std.nlower}{19}
\@namedef{rn@indep.IDRiD.TJDR-std.p}{$<$ 0.001}
\@namedef{rn@indep.IDRiD.TJDR-std.pct}{$-$33}
\@namedef{rn@indep.IDRiD.TJDR-std.rho}{0.51}
\@namedef{rn@indep.IDRiD.TJDR-uwf.drop}{0.42}
\@namedef{rn@indep.IDRiD.TJDR-uwf.first.intrank}{14}
\@namedef{rn@indep.IDRiD.TJDR-uwf.first.model}{MedSAM (full FT)}
\@namedef{rn@indep.IDRiD.TJDR-uwf.first.value}{0.295}
\@namedef{rn@indep.IDRiD.TJDR-uwf.gain}{-0.42}
\@namedef{rn@indep.IDRiD.TJDR-uwf.mean}{0.209}
\@namedef{rn@indep.IDRiD.TJDR-uwf.nlower}{19}
\@namedef{rn@indep.IDRiD.TJDR-uwf.p}{$<$ 0.001}
\@namedef{rn@indep.IDRiD.TJDR-uwf.pct}{$-$67}
\@namedef{rn@indep.IDRiD.TJDR-uwf.rho}{0.35}
\@namedef{rn@indep.IDRiD.fine.first.model}{SAMed}
\@namedef{rn@indep.IDRiD.fine.first.value}{0.394}
\@namedef{rn@indep.IDRiD.fine.hrnet}{0.367}
\@namedef{rn@indep.IDRiD.int.mean}{0.630}
\@namedef{rn@indep.IDRiD.ndrop}{19}
\@namedef{rn@indep.IDRiD.rho}{0.65}
\@namedef{rn@indep.IDRiD.sametop}{the same model}
\@namedef{rn@indep.IDRiD.top.model}{SAMed}
\@namedef{rn@indep.IDRiD.top.value}{0.39}
\@namedef{rn@indep.pairs.DiaRetDB1.ddrbetter}{13}
\@namedef{rn@indep.pairs.FGADR.ddrbetter}{14}
\@namedef{rn@indep.pairs.MAPLES-DR.ddrbetter}{13}
\@namedef{rn@indep.pairs.Retinal-Lesions.ddrbetter}{15}
\@namedef{rn@indep.pairs.TJDR-std.ddrbetter}{15}
\@namedef{rn@indep.pairs.TJDR-uwf.ddrbetter}{4}
\@namedef{rn@indep.pairs.ddrbetter}{74}
\@namedef{rn@indep.pairs.n}{114}
\@namedef{rn@indep.pairs.std.hi}{15}
\@namedef{rn@indep.pairs.std.lo}{13}
\@namedef{rn@lit.base.copied1}{22}
\@namedef{rn@lit.base.copied2}{26}
\@namedef{rn@lit.base.n}{1,538}
\@namedef{rn@lit.base.rerun}{28}
\@namedef{rn@lit.base.unverified}{23}
\@namedef{rn@lit.code.yes}{8}
\@namedef{rn@lit.ddr.n}{24}
\@namedef{rn@lit.ddr.notstated}{2}
\@namedef{rn@lit.ddr.official}{20}
\@namedef{rn@lit.ddr.own}{2}
\@namedef{rn@lit.deeplab.hi}{0.67}
\@namedef{rn@lit.deeplab.lo}{0.48}
\@namedef{rn@lit.deeplab.n}{11}
\@namedef{rn@lit.fov.notstated}{30}
\@namedef{rn@lit.hed.hi}{0.64}
\@namedef{rn@lit.hed.lo}{0.37}
\@namedef{rn@lit.hed.n}{7}
\@namedef{rn@lit.idrid.n}{30}
\@namedef{rn@lit.idrid.notstated}{3}
\@namedef{rn@lit.idrid.official}{27}
\@namedef{rn@lit.idrid.own}{0}
\@namedef{rn@lit.impl.exact}{2}
\@namedef{rn@lit.impl.fixed}{6}
\@namedef{rn@lit.impl.n}{27}
\@namedef{rn@lit.impl.notstated}{18}
\@namedef{rn@lit.impl.perimage}{1}
\@namedef{rn@lit.lseg.hi}{0.66}
\@namedef{rn@lit.lseg.lo}{0.40}
\@namedef{rn@lit.lseg.n}{14}
\@namedef{rn@lit.m2mrf.hi}{0.67}
\@namedef{rn@lit.m2mrf.lo}{0.58}
\@namedef{rn@lit.m2mrf.n}{6}
\@namedef{rn@lit.n}{30}
\@namedef{rn@lit.n.methods}{29}
\@namedef{rn@lit.pooled.stated}{5}
\@namedef{rn@lit.published.multi}{7}
\@namedef{rn@lit.rerun.all}{5}
\@namedef{rn@lit.rerun.copied}{3}
\@namedef{rn@lit.rerun.notstated}{10}
\@namedef{rn@lit.rerun.some}{12}
\@namedef{rn@lit.unet.hi}{0.65}
\@namedef{rn@lit.unet.lo}{0.35}
\@namedef{rn@lit.unet.n}{6}
\@namedef{rn@n.images}{5,121}
\@namedef{rn@n.independent}{6}
\@namedef{rn@n.main}{19}
\@namedef{rn@n.models}{39}
\@namedef{rn@n.ranked}{19}
\@namedef{rn@n.sources}{8}
\@namedef{rn@n.tier1}{22}
\@namedef{rn@n.tier2}{6}
\@namedef{rn@n.tier3}{11}
\@namedef{rn@params.nnunet}{59.2}
\@namedef{rn@params.resenc}{105.9}
\@namedef{rn@params.resenc.DDR}{140.0}
\@namedef{rn@params.resenc.IDRiD}{105.9}
\@namedef{rn@recall.DDR.large.best}{0.70}
\@namedef{rn@recall.DDR.large.median}{0.66}
\@namedef{rn@recall.DDR.small.best}{0.31}
\@namedef{rn@recall.DDR.small.median}{0.25}
\@namedef{rn@ref.EX.delta.mean}{-0.001}
\@namedef{rn@ref.EX.delta.min}{-0.002}
\@namedef{rn@ref.HE.delta.mean}{-0.000}
\@namedef{rn@ref.HE.delta.min}{-0.001}
\@namedef{rn@ref.MA.delta.mean}{-0.002}
\@namedef{rn@ref.MA.delta.min}{-0.004}
\@namedef{rn@ref.SE.delta.mean}{-0.000}
\@namedef{rn@ref.SE.delta.min}{-0.001}
\@namedef{rn@ref.delta.max}{0.000}
\@namedef{rn@ref.delta.mean}{-0.001}
\@namedef{rn@ref.delta.min}{-0.002}
\@namedef{rn@ref.geometry.absmax}{0.006}
\@namedef{rn@ref.geometry.tau}{0.95}
\@namedef{rn@ref.n}{19}
\@namedef{rn@ref.tau}{1.00}
\@namedef{rn@tau.published}{-0.11}
\@namedef{rn@tau.published.n}{10}
\@namedef{rn@tau.published.p}{0.73}
\@namedef{rn@vanish.DDR.1440.EX}{0.1}
\@namedef{rn@vanish.DDR.1440.MA}{0.0}
\@namedef{rn@vanish.DDR.224.EX}{35.3}
\@namedef{rn@vanish.DDR.224.MA}{34.3}
\@namedef{rn@vanish.DDR.384.EX}{16.9}
\@namedef{rn@vanish.DDR.384.MA}{11.5}
\@namedef{rn@vanish.IDRiD.1440.EX}{0.0}
\@namedef{rn@vanish.IDRiD.1440.MA}{0.0}
\@namedef{rn@vanish.IDRiD.224.EX}{25.8}
\@namedef{rn@vanish.IDRiD.224.MA}{16.0}
\@namedef{rn@vanish.IDRiD.384.EX}{5.7}
\@namedef{rn@vanish.IDRiD.384.MA}{1.1}
\makeatother


% --- text whose wording depends on simulated numbers
\ifretifinal
  \newenvironment{simtext}{}{}
\else
  \newtcolorbox{simtextbox}{enhanced,breakable,frame hidden,colback=white,boxrule=0pt,borderline west={1.2pt}{0pt}{phcolor},
    left=6pt,right=0pt,top=0pt,bottom=0pt,boxsep=0pt,before skip=4pt,after skip=4pt,
    before upper={\setlength{\parindent}{1.2em}\noindent}}
  \newenvironment{simtext}{\begin{simtextbox}}{\end{simtextbox}}
\fi

% --- headings in the template's style
\titleformat{\section}{\sffamily\bfseries\Large}{}{0pt}{}
\titleformat{\subsection}{\sffamily\bfseries\normalsize}{}{0pt}{}
\titleformat{\subsubsection}[runin]{\sffamily\bfseries\small}{}{0pt}{}[.]
\titlespacing*{\section}{0pt}{14pt}{4pt}
\titlespacing*{\subsection}{0pt}{10pt}{2pt}
\titlespacing*{\subsubsection}{0pt}{6pt}{4pt}
\setlength{\parskip}{0pt}
\setlength{\parindent}{1.2em}
\renewcommand{\figurename}{Fig.}
\captionsetup[table]{name=Table}

\newcommand{\figfile}[1]{figures/out/#1}
\newcommand{\widefig}[1]{\includegraphics[width=\textwidth,height=0.84\textheight,keepaspectratio]{\figfile{#1}}}
\newcommand{\mAUPR}{mAUPR}
\newcommand{\Bstd}{B-std}

\begin{document}
\linenumbers

\begin{flushleft}
{\sffamily\small\color{inkmuted} Analysis}\\[6pt]
{\sffamily\bfseries\LARGE Assessment of computational methods for diabetic retinopathy lesion segmentation in fundus images\par}
\vspace{8pt}
{\sffamily\normalsize Author One$^{1}$, Author Two$^{1}$, \ldots, Corresponding Author$^{1,\ast}$\par}
\vspace{3pt}
{\sffamily\footnotesize $^{1}$Affiliation. $^{\ast}$e-mail: corresponding@author\par}
\end{flushleft}

\ifretifinal\else
\begin{tcolorbox}[colback=headgray,colframe=rulegray,boxrule=0.4pt,arc=2pt,left=6pt,right=6pt,top=4pt,bottom=4pt,fontupper=\sffamily\footnotesize]
\textbf{Draft with simulated numbers.} Numbers printed in {\color{phcolor}this colour} are produced by \texttt{scripts/build\_placeholder.py} and are not results; text beside a vertical bar in the same colour was written against them and must be reread once the runs finish. Numbers in black are already measured or are frozen properties of the data and the protocol. Figure panels built on simulated rows carry a ``simulated'' tag. Set \texttt{\textbackslash retifinaltrue} to remove all three markers.
\end{tcolorbox}
\fi

% ------------------------------------------------------------------ Abstract
\begin{simtext}
\noindent{\sffamily\bfseries\small
Pixel-level segmentation of microaneurysms, hemorrhages, hard exudates and soft exudates makes automated grading of diabetic retinopathy interpretable, and dozens of methods now report results on the same few public datasets. Whether those numbers can be compared is unclear: splits, input resolution and metric implementations differ, and many baselines are copied from earlier papers. Here we assessed \rn{n.models} model variants on \rn{n.sources} test sets from six datasets (\rn{n.images} images) under one protocol. Of \rn{funnel.dr} fundus-specific methods, \rn{funnel.train} came with training code and \rn{funnel.full} with everything needed to rerun them; a widely reused 11-threshold implementation of the area under the precision--recall curve scored a released checkpoint at \rn{e1.m2mrf.IDRiD.fn} where the exact value was \rn{e1.m2mrf.IDRiD.exact}. After unified retraining of \rn{n.main} models, \rn{cd.first.model} had the best average rank, differences among the top \rn{cd.ntop} models were within the critical difference, and performance outside the training source depended on the annotation style and field of view of the test images: it held on pixel-labelled standard-field sources and fell by a third to a half on coarsely annotated, region-labelled and ultra-widefield images. Annotation granularity, input resolution and lesion size explained more variance than architecture. We release code, splits and probability maps as a benchmark.}
\end{simtext}

\vspace{8pt}

% -------------------------------------------------------------- Introduction
\noindent Diabetic retinopathy (DR) is graded from lesions that are visible in colour fundus photographs: microaneurysms (MA), hemorrhages (HE), hard exudates (EX) and soft exudates or cotton wool spots (SE). A model that outlines these lesions pixel by pixel shows a grader why an image received its grade, supports screening programmes in which the decision must be auditable, and provides quantitative lesion burden for follow-up\cite{sebastian2023drlesionsurvey,tahir2024microaneurysmreview}. Public datasets with pixel-level labels, first IDRiD\cite{porwal2018idrid,porwal2020idridchallenge} and DDR\cite{li2019ddr} and later FGADR\cite{zhou2021fgadr}, Retinal-Lesions\cite{wei2020retinallesionslesionnet} and TJDR\cite{mao2023tjdr}, have turned the task into a standard benchmark, and methods built for it now span multi-scale convolutional networks\cite{guo2019lseg,liu2023m2mrf,xu2024hacdrnet,he2022pmcnet}, relation and transformer designs\cite{huang2022rtnet,he2023h2former}, semi-supervised schemes\cite{zhou2019collaborativesemisupervised,ullah2023ssmdunet} and adaptations of segmentation foundation models\cite{kirillov2023sam,ravi2024sam2,ma2024medsam,li2024tpdrseg}.

Several properties of the task make the reported progress hard to read. (1) The datasets are small and heterogeneous. IDRiD has 81 annotated images, its official test set 27, and soft exudates appear in about half of them; other sources were annotated with ellipses and polygons by many graders, or as regions with confidence levels by four experts\cite{kauppi2007diaretdb1}, so that ``ground truth'' means different things from one source to the next. (2) The lesions are small. A microaneurysm spans about half a percent of the field-of-view diameter, and shrinking an image to the 224-pixel input of a common backbone removes \rn{vanish.DDR.224.MA}\% of the microaneurysms and \rn{vanish.DDR.224.EX}\% of the hard exudates from the DDR training labels before any model is trained. (3) Evaluation is not standardized. Papers differ in the split (official or their own), the resolution at which predictions are scored, whether pixels outside the field of view count, and how the area under the precision--recall curve (AUPR) is computed; a widely reused implementation samples the curve at 11 thresholds. (4) Baseline numbers travel between papers without being rerun, sometimes across several generations and across different splits. (5) Claims about foundation models mostly come from the papers that propose an adaptation, with zero-shot use, prompting from ground truth and fine-tuning reported side by side.

Earlier syntheses do not resolve these issues. Surveys collect the numbers that papers report\cite{sebastian2023drlesionsurvey,tahir2024microaneurysmreview}. The closest study, by Playout and Cheriet\cite{playout2024crossdatasetgeneralization}, trained U-Net-type baselines across five datasets and showed that annotation style governs cross-dataset transfer; it did not rerun fundus-specific methods, did not audit published numbers and did not include foundation models. Outside ophthalmology, standardized reassessments have repeatedly found that part of the reported progress disappears once baselines are configured with care\cite{isensee2021nnunet,isensee2024revisited} and that rankings on small test sets are fragile\cite{maierhein2018rankings}. A recent assessment of T cell receptor--epitope predictors\cite{lu2026tcrepitopenaturemethods} supplies the design we follow: evaluate the original models, retrain everything that can be retrained under one protocol, and judge both on independent data.

We therefore assessed methods for DR lesion segmentation along five lines. (1) Systematic collection: a fixed literature search and a ledger of 86 records, \rn{n.sources} test sets with frozen splits, and \rn{n.models} registered model variants with a reproducibility level for each. (2) Audit of published numbers: which split, which metric implementation and which provenance stand behind each reported value. (3) Original models: released checkpoints scored with the authors' metric and with ours, and the authors' pipelines rerun from scratch. (4) Unified retraining: one protocol for data, training amount, augmentation, checkpoint selection and evaluation, three seeds, and five test sets from sources that were never used for training. (5) Factors: annotation granularity, training source, training amount, input resolution, lesion size, loss and the way a foundation model is used, each varied alone. The study answers three questions: how many published numbers can be reproduced and where the gaps come from (Q1); which family of methods is better once conditions are equal, and whether the answer survives on independent data (Q2); and which properties of the data and the protocol decide performance (Q3).

\section{Results}

% ============================================================== R1
\subsection{Data collection and study design}
We designed a workflow that joins data collection, model collection and four kinds of evaluation (Fig.~\ref{fig:design}a). Seven test sets were prepared from six datasets, \rn{n.images} images in total (Fig.~\ref{fig:design}b and Table~\ref{tab:datasets}). IDRiD and DDR are the training sources of the main tables, because both have official splits and almost every published number refers to one of them. TJDR, split into its standard-field and 133$^\circ$ ultra-widefield parts, FGADR, Retinal-Lesions and DiaRetDB1 are independent sources: no model in the main tables was trained on them, and every checkpoint is applied to their test sets as it is. The sources differ in how lesions were drawn. Following the analysis of lesion count and lesion area by Playout and Cheriet\cite{playout2024crossdatasetgeneralization} we assigned each source to one of four granularity tiers, fine, mixed, coarse and region consensus, and we average results only within a tier.

A fixed search of three bibliographic databases identified 1,389 records; 764 remained after a Boolean gate and de-duplication, 112 entered the full-text queue, and screening of the priority list added 24 methods and 2 datasets to a snowballed corpus, for a ledger of 86 records (Fig.~\ref{fig:design}d). From it we registered \rn{n.models} model variants in five families and recorded for each a reproducibility level from R0 (no code) to R4 (training code, weights, splits and environment) (Fig.~\ref{fig:design}c and Table~\ref{tab:models}). \rn{n.tier1} variants form tier~1 and enter the main comparisons: general segmentation networks (U-Net\cite{ronneberger2015unet}, DeepLabv3\cite{chen2017deeplabv3,chen2018deeplabv3plus}, HRNet\cite{wang2020hrnet}, Swin-Unet\cite{cao2022swinunet}, FCT\cite{tragakis2023fct}, H2Former\cite{he2023h2former}), the self-configuring nnU-Net\cite{isensee2021nnunet}, the fundus-specific M2MRF\cite{liu2023m2mrf} and HACDR-Net\cite{xu2024hacdrnet}, and eight foundation models (SAM\cite{kirillov2023sam}, SAM~2\cite{ravi2024sam2}, SAM~3\cite{carion2025sam3}, MedSAM\cite{ma2024medsam}, SAMed\cite{zhang2023samed}, SAM-Adapter\cite{chen2023samadapter}, SAM2-UNet\cite{xiong2024sam2unet} and RETFound\cite{zhou2023retfound}). \rn{n.tier2} variants in tier~2 are added as time allows and \rn{n.tier3} in tier~3 are audited from their papers only.

Each model can appear under four evaluations that are never mixed in one column. The paper audit records the numbers as published and where they came from. Author weights are released checkpoints scored on official and independent test sets. Author recipes are the authors' pipelines rerun from scratch. Unified retraining applies one protocol, called \Bstd{} below, to the tier-1 models with complete training code. Two test scenarios are reported throughout, as in the template study\cite{lu2026tcrepitopenaturemethods}: the internal test, which is the official test set of the training source, and the independent test. The primary metric is the exact AUPR of each lesion class over all field-of-view pixels of a test set, averaged over the four classes (\mAUPR); Dice, intersection over union and lesion-level recall by lesion size are secondary. Finally, six factors are varied one at a time on one to five representative models (Fig.~\ref{fig:design}e).

\begin{figure}[p]
\centering
\widefig{fig1_design}
\caption{\textbf{Data collection and study design.} \textbf{a}, Workflow: the task, the seven test sources, the registered model variants, the four kinds of evaluation, and the two test scenarios with their metrics. \textbf{b}, Images per source and frozen split (train\,/\,validation\,/\,test). Colour gives the annotation granularity tier; shade separates the three parts of a split. \textbf{c}, Registered model variants by family and reproducibility level; the number of tier-1 variants is written after each bar. \textbf{d}, Literature search and screening. \textbf{e}, Factors varied one at a time, with the experiment code used in Methods.}
\label{fig:design}
\end{figure}

\begin{table*}[t]
\centering\tablefont
\caption{\textbf{Datasets, frozen splits and roles.} All eight sources carry pixel- or region-level labels for microaneurysms (MA), hemorrhages (HE), hard exudates (EX) and soft exudates (SE). Granularity decides which test sets may be averaged: averages are taken only within one tier. Split is train\,/\,validation\,/\,test in images; official test sets are never changed. ``Maps'' marks sources whose probability maps we can release.}
\label{tab:datasets}
\begin{adjustbox}{max width=\textwidth}
\begin{threeparttable}
\begin{tabular}{@{}l l l l l l c r c@{}}
\toprule
\colh{Source} & \colh{Ref.} & \colh{Images} & \colh{Resolution} & \colh{FOV} & \colh{Annotation} & \colh{Granularity} & \colh{Split} & \colh{Maps} \\
\midrule
\grouprow{9}{Training sources (internal test = official test set)}
IDRiD & \cite{porwal2018idrid} & 81 \inlinebar{0.14} & 4288$\times$2848 & 50$^\circ$ & pixel contours, expert review & \tierbadge{fine}{Fine} & 44 / 10 / 27 & \yes \\
DDR & \cite{li2019ddr} & 757 \inlinebar{1.32} & variable & 45$^\circ$ & pixel masks & \tierbadge{fine}{Fine} & 383 / 149 / 225 & \yes \\
\grouprow{9}{Independent test sources (never used to train a main-table model)}
TJDR, standard field & \cite{mao2023tjdr} & 257 \inlinebar{0.45} & 2048$\times$2048 & 35--50$^\circ$ & polygons, 2 graders + adjudicator & \tierbadge{fine}{Fine} & 172 / 30 / 55 & \yes \\
TJDR, ultra-widefield & \cite{mao2023tjdr} & 304 \inlinebar{0.53} & 3912$\times$3912 & 133$^\circ$ & polygons, 2 graders + adjudicator & \tierbadge{fine}{Fine} & 209 / 37 / 58 & \yes \\
MAPLES-DR & \cite{maplesdr2024scientificdata} & 198 \inlinebar{0.34} & 1440$\times$960–2304$\times$1536 & 45$^\circ$ & pixel masks, 7 retinologists; MA/HE/EX AI-assisted & \tierbadge{fine}{Fine} & \na\ / 0 / 60 & \yes \\
FGADR & \cite{zhou2021fgadr} & 1,842 \inlinebar{3.20} & 1280$\times$1280 & n.s. & pixel masks, 2 residents + senior & \tierbadge{mixed}{Mixed} & 1,096 / 193 / 553 & \no \\
Retinal-Lesions & \cite{wei2020retinallesionslesionnet} & 1,593 \inlinebar{2.77} & 896$\times$896 & n.s. & ellipses and polygons, $\geq$3 of 45 graders & \tierbadge{coarse}{Coarse} & 951 / 164 / 478 & \no \\
DiaRetDB1 & \cite{kauppi2007diaretdb1} & 89 \inlinebar{0.15} & 1500$\times$1152 & 50$^\circ$ & regions with confidence, 4 experts & \tierbadge{region}{Region} & \na\ / 28 / 61 & \tbd \\
\bottomrule
\end{tabular}
\begin{tablenotes}[flushleft]\scriptsize
\item FOV, field of view; n.s., not stated in the source paper. IDRiD: 10 of the 54 official training images are held out for validation in unified retraining; author-recipe retraining uses all 54. DiaRetDB1: ground truth is the mean confidence of four experts thresholded at 0.75; the 28 official training images are used only to choose thresholds. Retinal-Lesions: split by patient; grey-level 127 regions are ignored. MAPLES-DR: labels drawn at $1500\\times1500$ on the field-of-view crop of MESSIDOR images and placed back at the original resolution; the 138 training images are not used, the official 60 test images are; uncertain red (MA or HE) and bright (EX, SE or drusen) lesions are ignored for the two classes they could belong to. The FGADR licence forbids redistribution of derived data; maps for Retinal-Lesions are not released either.
\end{tablenotes}
\end{threeparttable}
\end{adjustbox}
\end{table*}

\begin{table*}[t]
\centering\tablefont
\caption{\textbf{Models entering the retraining tracks.} 38 variants are registered: 17 in tier~1 (main tables), 8 in tier~2 (if time allows) and 13 in tier~3 (paper audit only; Supplementary Table~2). Reproducibility level: R0, no code; R1, model definition only; R2, inference code and weights; R3, training code; R4, training code, weights, splits and environment. Tracks: A, author weights; B, unified retraining of the official code; C, re-implementation from the paper, reported in a separate column.}
\label{tab:models}
\begin{adjustbox}{max width=\textwidth}
\begin{threeparttable}
\begin{tabular}{@{}l c l c c c l l@{}}
\toprule
 & & & \multicolumn{3}{c}{\colh{Track}} & & \\
\cmidrule(lr){4-6}
\colh{Model} & \colh{Tier} & \colh{Reproducibility} & \colh{A} & \colh{B} & \colh{C} & \colh{Code status} & \colh{Use in this study} \\
\midrule
\grouprow{8}{General segmentation and self-configuring baselines}
\famchip{general}\; U-Net & 1 & \rlevel{3} {\scriptsize R3} & \no & \yes & \no & reference implementation & B-std \\
\famchip{general}\; DeepLabv3 & 1 & \rlevel{3} {\scriptsize R3} & \no & \yes & \no & torchvision, ResNet-101 & B-std \\
\famchip{general}\; HRNet-W48 & 1 & \rlevel{3} {\scriptsize R3} & \no & \yes & \no & official, W48 & B-std \\
\famchip{auto}\; nnU-Net & 1 & \rlevel{4} {\scriptsize R4} & \no & \yes & \no & official, 2D; default and ResEnc-M & own entry point \\
\famchip{general}\; Swin-Unet & 1 & \rlevel{3} {\scriptsize R3} & \no & \yes & \no & official & B-std, 224 patches \\
\famchip{general}\; FCT & 1 & \rlevel{3} {\scriptsize R3} & \no & \yes & \no & official PyTorch port & B-std, 384 patches \\
\famchip{general}\; H2Former & 1 & \rlevel{3} {\scriptsize R3} & \no & \yes & \no & official & deferred: exceeds 40\,GB at 960 \\
\famchip{general}\; UNet++ & 2 & \rlevel{3} {\scriptsize R3} & \no & \yes & \no & official & tier 2 \\
\famchip{general}\; UPerNet & 2 & \rlevel{3} {\scriptsize R3} & \no & \yes & \no & official & tier 2 \\
\grouprow{8}{Fundus-specific, multi-scale CNN}
\famchip{fundus}\; M2MRF & 1 & \rlevel{4} {\scriptsize R4} & \yes & \yes & \no & official; weights released & A, E1r, B-std \\
\famchip{fundus}\; HACDR-Net & 1 & \rlevel{3} {\scriptsize R3} & \no & \yes & \no & official; no weights & E1r, B-std \\
\famchip{fundus}\; WSRFNet & 2 & \rlevel{1} {\scriptsize R1} & \yes & \no & \yes & model and test code only & A if weights arrive; C \\
\famchip{fundus}\; L-Seg & 2 & \rlevel{1} {\scriptsize R1} & \no & \no & \yes & Caffe prototxt and loss layer & C \\
\famchip{fundus}\; SAA & 2 & \rlevel{0} {\scriptsize R0} & \no & \no & \yes & repository empty & C \\
\famchip{fundus}\; PMCNet & 2 & \rlevel{3} {\scriptsize R3} & \no & \yes & \no & official, TF\,1.13; no weights & tier 2 \\
\grouprow{8}{Fundus-specific, relation or transformer}
\famchip{fundus}\; RTNet & 2 & \rlevel{0} {\scriptsize R0} & \no & \no & \yes & no code; image + vessel input & C \\
\grouprow{8}{Foundation models}
\famchip{foundation}\; SAM & 1 & \rlevel{2} {\scriptsize R2} & \no & \yes & \no & official inference & zero-shot, oracle prompt \\
\famchip{foundation}\; SAM 2 & 1 & \rlevel{3} {\scriptsize R3} & \no & \yes & \no & official & zero-shot, oracle prompt \\
\famchip{foundation}\; SAM 3 & 1 & \rlevel{3} {\scriptsize R3} & \no & \yes & \no & official; gated weights & text prompt \\
\famchip{foundation}\; MedSAM & 1 & \rlevel{3} {\scriptsize R3} & \no & \yes & \no & official & oracle box, full fine-tuning \\
\famchip{foundation}\; SAMed & 1 & \rlevel{3} {\scriptsize R3} & \no & \yes & \no & official & PEFT (LoRA) \\
\famchip{foundation}\; SAM-Adapter & 1 & \rlevel{3} {\scriptsize R3} & \no & \yes & \no & official & PEFT \\
\famchip{foundation}\; SAM2-UNet & 1 & \rlevel{3} {\scriptsize R3} & \no & \yes & \no & official & PEFT \\
\famchip{foundation}\; RETFound + UPerNet & 1 & \rlevel{3} {\scriptsize R3} & \no & \yes & \no & official encoder; UPerNet head added & encoder + decoder \\
\famchip{foundation}\; PerSAM & 2 & \rlevel{2} {\scriptsize R2} & \no & \yes & \no & official inference & one-shot, tier 2 \\
\bottomrule
\end{tabular}
\begin{tablenotes}[flushleft]\scriptsize
\item \famchip{general} general segmentation \quad \famchip{auto} self-configuring \quad \famchip{fundus} fundus-specific \quad \famchip{foundation} foundation model. B-std, unified retraining protocol (Methods). PEFT, parameter-efficient fine-tuning. Status as of the code audit of 2--4 October 2026; author replies are due by 31 October 2026.
\end{tablenotes}
\end{threeparttable}
\end{adjustbox}
\end{table*}


% ============================================================== R2
\subsection{The numbers in the literature}
We first asked what stands behind the numbers that papers report. Of the \rn{funnel.dr} fundus-specific methods in the registry (17 dedicated architectures and 3 adaptations of foundation models to this task), \rn{funnel.code} have any author code, \rn{funnel.train} have training code and \rn{funnel.full} (M2MRF) comes with weights, splits and an environment description (Fig.~\ref{fig:literature}a). Authors of methods without code or weights were contacted with a fixed template; methods that could not be run stay in the denominator.

For every method paper in the ledger (\rn{lit.n.methods} method papers and one prior cross-dataset benchmark\cite{playout2024crossdatasetgeneralization}) we extracted the experimental setup and the comparison tables, \rn{lit.base.n} baseline AUPR values on IDRiD and DDR in all (Supplementary Table~3). The official split is the rule: \rn{lit.idrid.official} of the \rn{lit.idrid.n} papers that use IDRiD and \rn{lit.ddr.official} of the \rn{lit.ddr.n} that use DDR state it, and the others give no split or merge the official test images into training. What the papers leave out is how the metric was computed. Of \rn{lit.impl.n} papers whose main metric is AUPR, \rn{lit.impl.notstated} give no implementation, \rn{lit.impl.fixed} integrate over a fixed grid of 11 or 33 thresholds, which is the IDRiD challenge protocol carried on by the evaluation code released with L-Seg and M2MRF, \rn{lit.impl.exact} use the exact average precision and \rn{lit.impl.perimage} averages over images; none states that the evaluation is restricted to the field of view, and \rn{lit.pooled.stated} say whether pixels are pooled over the test set or scores are averaged over images (Fig.~\ref{fig:literature}b). Baseline numbers travel between papers. Of the baseline values, \rn{lit.base.rerun}\% were produced by the paper that reports them, \rn{lit.base.copied1}\% are identical to the value in the method's own paper, \rn{lit.base.copied2}\% are identical to a value that an earlier paper had already copied or produced, and \rn{lit.base.unverified}\% cite a source outside the ledger or none. Only \rn{lit.rerun.all} papers state that every baseline was rerun, and \rn{lit.rerun.notstated} say nothing about where their comparison numbers come from. As a consequence one method carries a range of values: L-Seg is reported on IDRiD by \rn{lit.lseg.n} papers with \mAUPR{} from \rn{lit.lseg.lo} to \rn{lit.lseg.hi}, DeepLabv3+ by \rn{lit.deeplab.n} papers from \rn{lit.deeplab.lo} to \rn{lit.deeplab.hi} and U-Net by \rn{lit.unet.n} papers from \rn{lit.unet.lo} to \rn{lit.unet.hi}. Copied values agree to the last digit; the spread comes from the papers that reran a method under their own settings (Fig.~\ref{fig:literature}c). The audit also surfaced copied tables whose lesion columns are transposed with respect to the source, a baseline table reproduced verbatim from an earlier paper of the same group without attribution, and a Dice value carried into an AUPR column (Supplementary Table~3).

The metric implementation alone moves the numbers. The evaluation code released with M2MRF, which later work reuses, computes AUPR from 11 thresholds; at threshold 0 every pixel with a non-zero probability is counted as positive, so the recall is 1 and the precision equals the prevalence, which lifts the area under the curve. On the released IDRiD checkpoint this function returns \rn{e1.m2mrf.IDRiD.fn}; after the probabilities are stored as 16-bit maps the same function returns \rn{e1.m2mrf.IDRiD.stored}, and the exact pooled AUPR inside the field of view is \rn{e1.m2mrf.IDRiD.exact}. On DDR the three values are \rn{e1.m2mrf.DDR.fn}, \rn{e1.m2mrf.DDR.stored} and \rn{e1.m2mrf.DDR.exact} (Fig.~\ref{fig:literature}d).
\begin{simtext}
Applied to the probability maps of all retrained models, the 11-threshold implementation raised \mAUPR{} on IDRiD by \rn{impl.mean} on average and by up to \rn{impl.max}, by an amount that depends on how each model distributes its low probabilities. \rn{impl.moved} of \rn{n.main} models changed rank (Kendall $\tau$ = \rn{impl.tau}; Fig.~\ref{fig:literature}e).
\end{simtext}

\begin{figure}[t]
\centering
\widefig{fig2_literature}
\caption{\textbf{The numbers in the literature.} \textbf{a}, From registered fundus-specific methods to a reproduced number; the first four stages follow the reproducibility levels of the registry. \textbf{b}, Share of method papers by split and by AUPR implementation (fixed thresholds, a grid of 11 or 33 operating points), and provenance of the baseline AUPR values on IDRiD and DDR in their comparison tables: produced by the reporting paper, identical to the method's own paper, identical to a value already carried by an earlier paper, or citing a source outside the ledger. \textbf{c}, \mAUPR{} on IDRiD reported for one method in different papers (dots; bar, range; $n$, number of papers; methods reported by at least three papers) and the value after unified retraining (diamond). \textbf{d}, The same probability maps scored with the exact AUPR and with the 11-threshold implementation; the two shaded rows are the released M2MRF checkpoints, for which the authors' function on the float output is also shown. \textbf{e}, Rank of the retrained models under the two implementations.}
\label{fig:literature}
\end{figure}

% ============================================================== R3
\subsection{Performance of the original models}
Released checkpoints were scored with three numbers each: the value printed in the paper, the authors' evaluation function, and our protocol (Fig.~\ref{fig:original}a and Table~\ref{tab:original}). For M2MRF the authors' function reproduced the repository values to the last printed digit on both datasets (\rn{e1.m2mrf.IDRiD.fn} and \rn{e1.m2mrf.DDR.fn}), which shows that moving the code to a current framework did not change its output; the paper reports the mean of three runs (\rn{e1.m2mrf.IDRiD.paper} and \rn{e1.m2mrf.DDR.paper}). Under our protocol the same checkpoints reach \rn{e1.m2mrf.IDRiD.exact} and \rn{e1.m2mrf.DDR.exact}, a gap of \rn{e1.m2mrf.IDRiD.gap} and \rn{e1.m2mrf.DDR.gap} that is almost entirely due to the threshold-0 term of the 11-threshold curve; restricting the evaluation to the field of view and keeping overlapping labels each change the third decimal only (Fig.~\ref{fig:original}b).

Rerunning the authors' recipe is the direct test of whether released code reproduces a reported number. A first rerun of M2MRF on IDRiD, with an effective batch of 1 instead of the authors' 4, reached \rn{e1r.first.fn} with the authors' metric and \rn{e1r.first.exact} with ours, well below the released checkpoint.
\begin{simtext}
With the authors' batch size, M2MRF came within \rn{e1r.M2MRF.IDRiD.delta} of the reported value on IDRiD over three seeds and within \rn{e1r.M2MRF.DDR.delta} on DDR in one run, and HACDR-Net, for which no weights are released, within \rn{e1r.HACDR-Net.IDRiD.delta} and \rn{e1r.HACDR-Net.DDR.delta} in one run per source (Fig.~\ref{fig:original}c). On the independent test sets the author weights lost most of their internal performance, with the same ordering of sources as for the retrained models below (Fig.~\ref{fig:original}d).
\end{simtext}

\begin{figure}[t]
\centering
\widefig{fig3_original}
\caption{\textbf{Performance of the original models.} \textbf{a}, \mAUPR{} of each released checkpoint on the official test set of its training source: reported in the paper, authors' function on the checkpoint, and this study's protocol. The number on the right is the difference between the last two. \textbf{b}, Decomposition of that difference for the two M2MRF checkpoints. \textbf{c}, Retraining with the authors' recipe: reported value (line) and individual seeds (dots); the cross is the first rerun at an effective batch of 1. \textbf{d}, Per-class AUPR of the author weights on all seven test sets under this study's protocol. Outlined cells are internal tests.}
\label{fig:original}
\end{figure}

\begin{table*}[t]
\centering\tablefont
\caption{\textbf{Original models: three numbers per released checkpoint, and retraining with the authors' recipe.} \textbf{a}, Each checkpoint is scored on the official test set of its training source with the number printed in the paper, the authors' evaluation function (11 thresholds, whole image, overwritten labels) on the float output, the same function on the stored 16-bit maps, and this study's protocol (exact pooled AUPR inside the field of view, overlapping labels). \textbf{b}, The authors' pipeline rerun from scratch; mean $\pm$ s.d. over seeds, last iteration, no checkpoint selection.}
\label{tab:original}
\begin{adjustbox}{max width=\textwidth}
\begin{threeparttable}
\begin{tabular}{@{}l l c c c c c c c c c c@{}}
\toprule
\multicolumn{12}{@{}l}{\textbf{a}\quad Author weights (E1)}\\
\addlinespace[2pt]
 & & \multicolumn{3}{c}{\colh{mAUPR, author metric}} & \multicolumn{5}{c}{\colh{AUPR, this protocol}} & & \\
\cmidrule(lr){3-5}\cmidrule(lr){6-10}
\colh{Model} & \colh{Source} & Paper & Function & Stored & \colh{mAUPR} & MA & HE & EX & SE & \colh{$\Delta$ metric} & \colh{Indep.\ fine} \\
\midrule
M2MRF & IDRiD & .672 & .676 & .645 & .639 & \cellcolor[HTML]{C0E4D8}.408 & \cellcolor[HTML]{88CEC9}.653 & \cellcolor[HTML]{5BB9C0}.832 & \cellcolor[HTML]{85CDC8}.664 & \down{.036} & \ph{.240} \\
M2MRF & DDR & .489 & .492 & .426 & .401 & \na & \na & \na & \na & \down{.091} & \ph{.276} \\
\addlinespace[2pt]
HRDecoder & IDRiD & \ph{.713} & \ph{.713} & \ph{.687} & \ph{.667} & \cellcolor[HTML]{B9E1D5}\ph{.449} & \cellcolor[HTML]{7CC9C6}\ph{.696} & \cellcolor[HTML]{59B7BF}\ph{.849} & \cellcolor[HTML]{83CCC8}\ph{.673} & \down{\ph{.047}} & \ph{.269} \\
HRDecoder & DDR & \ph{.492} & \ph{.492} & \ph{.449} & \ph{.429} & \cellcolor[HTML]{DFF0E5}\ph{.240} & \cellcolor[HTML]{AFDDD1}\ph{.500} & \cellcolor[HTML]{9CD6CC}\ph{.577} & \cellcolor[HTML]{C2E5D8}\ph{.400} & \down{\ph{.063}} & \ph{.310} \\
\addlinespace[2pt]
WSRFNet & IDRiD & \ph{.690} & \ph{.676} & \ph{.665} & \ph{.645} & \cellcolor[HTML]{BEE3D7}\ph{.424} & \cellcolor[HTML]{88CEC9}\ph{.650} & \cellcolor[HTML]{5AB8BF}\ph{.839} & \cellcolor[HTML]{85CDC8}\ph{.665} & \down{\ph{.032}} & \ph{.228} \\
WSRFNet & DDR & \ph{.500} & \ph{.494} & \ph{.425} & \ph{.405} & \cellcolor[HTML]{E1F1E6}\ph{.225} & \cellcolor[HTML]{B7E0D4}\ph{.465} & \cellcolor[HTML]{9ED7CD}\ph{.569} & \cellcolor[HTML]{C8E7DB}\ph{.364} & \down{\ph{.089}} & \ph{.254} \\
\bottomrule
\end{tabular}

\vspace{6pt}
\begin{tabular}{@{}l l c c c c c@{}}
\toprule
\multicolumn{7}{@{}l}{\textbf{b}\quad Author recipe, retrained (E1r)}\\
\addlinespace[2pt]
\colh{Model} & \colh{Source} & \colh{Paper} & \colh{Retrained, author metric} & \colh{Seeds} & \colh{$\Delta$ to paper} & \colh{Retrained, this protocol} \\
\midrule
M2MRF & IDRiD & .672 & \ph{.612}\sd{\ph{.154}} & \ph{3} & \down{\ph{.060}} & \ph{.549}\sd{\ph{.198}} \\
\quad first run, batch 1 & IDRiD & .672 & .576 & 1 & \down{.096} & .493 \\
M2MRF & DDR & .489 & \ph{.454}\sd{\ph{.002}} & \ph{3} & \down{\ph{.036}} & \ph{.416}\sd{\ph{.008}} \\
HACDR-Net & IDRiD & \ph{.700} & \ph{.596}\sd{\ph{.011}} & \ph{3} & \down{\ph{.104}} & \ph{.563}\sd{\ph{.016}} \\
HACDR-Net & DDR & \ph{.520} & \ph{.426}\sd{\ph{.006}} & \ph{3} & \down{\ph{.094}} & \ph{.389}\sd{\ph{.004}} \\
\bottomrule
\end{tabular}
\begin{tablenotes}[flushleft]\scriptsize
\item ``$\Delta$ metric'' is this protocol minus the authors' function on the same checkpoint. ``Indep.\ fine'' is the mean mAUPR on TJDR standard-field and ultra-widefield test sets. The M2MRF paper reports the mean of three runs (67.24 on IDRiD, 48.94 on DDR), the repository the best run (67.55, 49.20). HRDecoder: the released checkpoints (repository table 71.34, 49.24) through the authors' test pipeline (resize to $2880\times1920$ with sliding $1920$ windows on IDRiD, $2048\times2048$ whole image on DDR). The first M2MRF rerun used an effective batch of 1 instead of the authors' 4.
\phnote
\end{tablenotes}
\end{threeparttable}
\end{adjustbox}
\end{table*}


% ============================================================== R4
\subsection{A unified protocol to retrain the models}
Source recipes cannot serve as a common ground for retraining. The recipes of the general-purpose networks were written for natural images or abdominal organs; moved to IDRiD unchanged they differ about 67-fold in the number of training images seen, and those that resize the whole image to a small fixed input discard a large part of the small lesions (Fig.~\ref{fig:scale}b). We therefore fixed everything that is not part of a method: data and splits, a canvas of 1,440 pixels across the field of view, the training amount counted in pixels, augmentation, selection of the checkpoint by validation \mAUPR, inference geometry and evaluation. Each model keeps its loss, optimizer, learning rate, weight decay, schedule shape and effective batch (Methods and Table~\ref{tab:recipes}). The training amount was frozen from the validation curves of two models; no test metric was used. We retrained \rn{n.main} models on IDRiD and on DDR, with three seeds for the seven networks that have a source recipe and for the self-configuring baseline, and one seed for the adapted foundation models (two of them, which also serve as representative models below, with three seeds where noted); the rank tests use seed 0 of every model. The source recipes, for the models for which they completed, are reported in Supplementary Table~4.

\subsection{Performance of retrained models on internal tests}
\begin{simtext}
On the internal tests (Fig.~\ref{fig:retrained}a and Table~\ref{tab:main}), \rn{bstd.IDRiD.best.model} reached the highest \mAUPR{} on IDRiD (\rn{bstd.IDRiD.best.value}), followed by \rn{bstd.IDRiD.second.model} (\rn{bstd.IDRiD.second.value}); on DDR the first two were \rn{bstd.DDR.best.model} (\rn{bstd.DDR.best.value}) and \rn{bstd.DDR.second.model} (\rn{bstd.DDR.second.value}). The difference between the first two models was \rn{bstd.IDRiD.gap12} on IDRiD and \rn{bstd.DDR.gap12} on DDR, against a median seed standard deviation of \rn{bstd.IDRiD.sd.median} and \rn{bstd.DDR.sd.median}. IDRiD has 27 test images, and its bootstrap intervals are correspondingly wide; we use it to connect to the literature and base comparisons between architectures on DDR and on the independent tests.

Lesion classes differed more than models. On DDR the median AUPR over models was \rn{bstd.DDR.EX.median} for hard exudates and \rn{bstd.DDR.MA.median} for microaneurysms, and the best model reached \rn{bstd.DDR.MA.best} for microaneurysms. Networks with a coarse output stride or 16-pixel tokens lost most on this class. One failure was specific to a loss rather than to an architecture: the reference U-Net, trained on patches with a per-image Dice term smoothed by 1, drove its rare-class channels towards zero (soft exudate AUPR \rn{bstd.DDR.unet.SE} on DDR), because 83\% of DDR training patches contain no soft exudate and an empty patch rewards an all-zero output (Methods).
\end{simtext}

\begin{table*}[t]
\centering\tablefont
\caption{\textbf{Unified retraining (B-std): internal tests.} Every model is trained on one source under the shared protocol and scored on that source's official test set. Per-class values are exact pooled AUPR inside the field of view, mean of three seeds; mAUPR is the mean of the four classes $\pm$ s.d.\ over seeds; Dice is pooled, at the per-class threshold chosen on the validation set. Cell shading follows AUPR. \best{Bold}, best in column; \second{underlined}, second.}
\label{tab:main}
\begin{adjustbox}{max width=\textwidth}
\begin{threeparttable}
\begin{tabular}{@{}l r cccc c c @{\hspace{10pt}} cccc c c@{}}
\toprule
 & & \multicolumn{6}{c}{\colh{Trained and tested on IDRiD} {\scriptsize(44 / 10 / 27)}} & \multicolumn{6}{c}{\colh{Trained and tested on DDR} {\scriptsize(383 / 149 / 225)}} \\
\cmidrule(lr){3-8}\cmidrule(l){9-14}
\colh{Model} & \colh{Params} & MA & HE & EX & SE & \colh{mAUPR} & \colh{Dice} & MA & HE & EX & SE & \colh{mAUPR} & \colh{Dice} \\
\midrule
\grouprow{14}{\famchip{general}\; General segmentation}
U-Net & 7.8 & \cellcolor[HTML]{CAE8DC}\ph{.351} & \cellcolor[HTML]{99D5CC}\ph{.588} & \cellcolor[HTML]{66C1C2}\ph{.779} & \cellcolor[HTML]{E1F1E6}\ph{.222} & \ph{.485}\sd{\ph{.013}} & \ph{.478} & \cellcolor[HTML]{EEF5ED}\ph{.080} & \cellcolor[HTML]{E1F1E6}\ph{.225} & \cellcolor[HTML]{AFDDD1}\ph{.501} & \cellcolor[HTML]{EEF5ED}\ph{.080} & \ph{.222}\sd{\ph{.009}} & \ph{.271} \\
DeepLabv3 & \ph{60.9} & \cellcolor[HTML]{DEEFE4}\ph{.244} & \cellcolor[HTML]{A5D9CE}\ph{.539} & \cellcolor[HTML]{6FC4C4}\ph{.748} & \cellcolor[HTML]{A1D8CE}\ph{.556} & \ph{.522}\sd{\ph{.023}} & \ph{.506} & \cellcolor[HTML]{ECF4EB}\ph{.107} & \cellcolor[HTML]{CAE8DC}\ph{.356} & \cellcolor[HTML]{B8E1D4}\ph{.458} & \cellcolor[HTML]{D9EDE2}\ph{.270} & \ph{.298}\sd{\ph{.001}} & \ph{.328} \\
HRNet-W48 & \ph{65.8} & \cellcolor[HTML]{BBE2D6}\ph{.436} & \cellcolor[HTML]{88CEC9}\ph{.652} & \cellcolor[HTML]{5BB9C0}\ph{.833} & \cellcolor[HTML]{85CDC8}\ph{.663} & \ph{.646}\sd{\ph{.011}} & \ph{.603} & \cellcolor[HTML]{E2F1E6}\ph{.219} & \cellcolor[HTML]{B5E0D3}\ph{.469} & \cellcolor[HTML]{9CD6CC}\ph{.574} & \cellcolor[HTML]{C6E6DA}\ph{.373} & \ph{.409}\sd{\ph{.013}} & \ph{.416} \\
Swin-Unet & \ph{27.2} & \cellcolor[HTML]{DEEFE4}\ph{.242} & \cellcolor[HTML]{ACDCD0}\ph{.517} & \cellcolor[HTML]{71C5C4}\ph{.739} & \cellcolor[HTML]{A5D9CE}\ph{.543} & \ph{.511}\sd{\ph{.019}} & \ph{.500} & \cellcolor[HTML]{ECF5EC}\ph{.102} & \cellcolor[HTML]{D0EADE}\ph{.321} & \cellcolor[HTML]{C0E4D8}\ph{.407} & \cellcolor[HTML]{DCEFE4}\ph{.250} & \ph{.270}\sd{\ph{.006}} & \ph{.308} \\
FCT & \ph{51.4} & \cellcolor[HTML]{DCEFE4}\ph{.251} & \cellcolor[HTML]{AEDDD0}\ph{.508} & \cellcolor[HTML]{77C7C5}\ph{.721} & \cellcolor[HTML]{ACDCD0}\ph{.516} & \ph{.499}\sd{\ph{.012}} & \ph{.485} & \cellcolor[HTML]{ECF5EC}\ph{.105} & \cellcolor[HTML]{D3EBE0}\ph{.307} & \cellcolor[HTML]{C2E5D8}\ph{.400} & \cellcolor[HTML]{E0F0E5}\ph{.232} & \ph{.261}\sd{\ph{.015}} & \ph{.303} \\
\grouprow{14}{\famchip{auto}\; Self-configuring}
nnU-Net & \ph{33.4} & \cellcolor[HTML]{B3DFD2}\ph{\best{.479}} & \cellcolor[HTML]{81CBC7}\ph{\second{.679}} & \cellcolor[HTML]{58B6BF}\ph{\best{.856}} & \cellcolor[HTML]{7CC9C6}\ph{\best{.698}} & \ph{\best{.678}}\sd{\ph{.020}} & \ph{\second{.629}} & \cellcolor[HTML]{DEEFE4}\ph{\second{.248}} & \cellcolor[HTML]{B4DFD3}\ph{\second{.477}} & \cellcolor[HTML]{9CD6CC}\ph{\second{.579}} & \cellcolor[HTML]{C3E5D9}\ph{\second{.397}} & \ph{\second{.425}}\sd{\ph{.015}} & \ph{\second{.430}} \\
nnU-Net ResEnc-M & \ph{57.5} & \cellcolor[HTML]{B4DFD3}\ph{\second{.473}} & \cellcolor[HTML]{80CBC7}\ph{\best{.688}} & \cellcolor[HTML]{58B6BF}\ph{\second{.854}} & \cellcolor[HTML]{7ECAC6}\ph{\second{.692}} & \ph{\second{.677}}\sd{\ph{.006}} & \ph{\best{.629}} & \cellcolor[HTML]{DBEEE3}\ph{\best{.256}} & \cellcolor[HTML]{B0DED1}\ph{\best{.497}} & \cellcolor[HTML]{97D4CB}\ph{\best{.597}} & \cellcolor[HTML]{C2E5D8}\ph{\best{.402}} & \ph{\best{.438}}\sd{\ph{.022}} & \ph{\best{.448}} \\
\grouprow{14}{\famchip{fundus}\; Fundus-specific}
M2MRF & 70.3 & \cellcolor[HTML]{BDE3D6}\ph{.427} & \cellcolor[HTML]{8ED0CA}\ph{.628} & \cellcolor[HTML]{5DBAC0}\ph{.823} & \cellcolor[HTML]{88CEC9}\ph{.649} & \ph{.632}\sd{\ph{.007}} & \ph{.594} & \cellcolor[HTML]{E0F0E5}\ph{.234} & \cellcolor[HTML]{B5E0D3}\ph{.470} & \cellcolor[HTML]{9ED7CD}\ph{.568} & \cellcolor[HTML]{C5E6DA}\ph{.380} & \ph{.413}\sd{\ph{.022}} & \ph{.420} \\
HACDR-Net & \ph{49.6} & \cellcolor[HTML]{CFEADE}\ph{.323} & \cellcolor[HTML]{9CD6CC}\ph{.577} & \cellcolor[HTML]{65C0C2}\ph{.783} & \cellcolor[HTML]{99D5CC}\ph{.591} & \ph{.569}\sd{\ph{.019}} & \ph{.543} & \cellcolor[HTML]{E6F3E8}\ph{.169} & \cellcolor[HTML]{C0E4D8}\ph{.407} & \cellcolor[HTML]{AEDDD0}\ph{.511} & \cellcolor[HTML]{D0EADE}\ph{.320} & \ph{.352}\sd{\ph{.017}} & \ph{.377} \\
\grouprow{14}{\famchip{foundation}\; Foundation models, adapted}
SAMed & \ph{91.2} & \cellcolor[HTML]{DEEFE4}\ph{.243} & \cellcolor[HTML]{A5D9CE}\ph{.544} & \cellcolor[HTML]{6AC2C3}\ph{.766} & \cellcolor[HTML]{A0D7CD}\ph{.565} & \ph{.530}\sd{\ph{.009}} & \ph{.512} & \cellcolor[HTML]{EDF5EC}\ph{.096} & \cellcolor[HTML]{CDE9DD}\ph{.340} & \cellcolor[HTML]{BAE2D5}\ph{.445} & \cellcolor[HTML]{DAEEE3}\ph{.267} & \ph{.287}\sd{\ph{.004}} & \ph{.327} \\
SAM-Adapter & 641.3 & \cellcolor[HTML]{E2F1E6}\ph{.216} & \cellcolor[HTML]{AADCCF}\ph{.523} & \cellcolor[HTML]{75C6C5}\ph{.724} & \cellcolor[HTML]{A7DACF}\ph{.533} & \ph{.499}\sd{\ph{.019}} & \ph{.491} & \cellcolor[HTML]{EEF5ED}\ph{.085} & \cellcolor[HTML]{D0EADE}\ph{.322} & \cellcolor[HTML]{C2E5D8}\ph{.401} & \cellcolor[HTML]{DFF0E5}\ph{.239} & \ph{.262}\sd{\ph{.009}} & \ph{.306} \\
SAM2-UNet & 216.5 & \cellcolor[HTML]{C9E7DB}\ph{.361} & \cellcolor[HTML]{8ED0CA}\ph{.632} & \cellcolor[HTML]{5EBBC0}\ph{.817} & \cellcolor[HTML]{8CD0C9}\ph{.640} & \ph{.613}\sd{\ph{.026}} & \ph{.578} & \cellcolor[HTML]{E5F2E7}\ph{.194} & \cellcolor[HTML]{BAE2D5}\ph{.442} & \cellcolor[HTML]{A3D9CE}\ph{.551} & \cellcolor[HTML]{C8E7DB}\ph{.366} & \ph{.388}\sd{\ph{.010}} & \ph{.403} \\
MedSAM (full FT) & 93.7 & \cellcolor[HTML]{E3F1E6}\ph{.211} & \cellcolor[HTML]{A7DACF}\ph{.535} & \cellcolor[HTML]{6CC3C3}\ph{.757} & \cellcolor[HTML]{A5D9CE}\ph{.540} & \ph{.511}\sd{\ph{.022}} & \ph{.493} & \cellcolor[HTML]{EDF5EC}\ph{.089} & \cellcolor[HTML]{CEE9DD}\ph{.330} & \cellcolor[HTML]{BDE3D6}\ph{.428} & \cellcolor[HTML]{DBEEE3}\ph{.257} & \ph{.276}\sd{\ph{.001}} & \ph{.311} \\
RETFound + UPerNet & 352.6 & \cellcolor[HTML]{D4ECE0}\ph{.302} & \cellcolor[HTML]{95D3CB}\ph{.604} & \cellcolor[HTML]{63BFC1}\ph{.795} & \cellcolor[HTML]{91D2CA}\ph{.616} & \ph{.579}\sd{\ph{.013}} & \ph{.555} & \cellcolor[HTML]{E9F4EA}\ph{.134} & \cellcolor[HTML]{C3E5D9}\ph{.394} & \cellcolor[HTML]{AFDDD1}\ph{.499} & \cellcolor[HTML]{D1EBDF}\ph{.311} & \ph{.334}\sd{\ph{.024}} & \ph{.358} \\
\bottomrule
\end{tabular}
\begin{tablenotes}[flushleft]\scriptsize
\item Params in millions. nnU-Net uses its own preprocessing and planning with the same splits, test images and evaluation. H2Former is registered in tier~1 but deferred: at its 960-pixel input one training step exceeds the 40\,GB of the cards used here.
\phnote
\end{tablenotes}
\end{threeparttable}
\end{adjustbox}
\end{table*}


\subsection{Performance of retrained models on independent tests}
\begin{simtext}
Applied to sources that were never used for training, every model lost performance (Fig.~\ref{fig:retrained}b,c and Table~\ref{tab:independent}). For models trained on DDR the mean loss in \mAUPR{} was \rn{indep.DDR.TJDR-std.drop} on the standard-field TJDR test set (Holm-adjusted $P$ \rn{indep.DDR.TJDR-std.p}), \rn{indep.DDR.TJDR-uwf.drop} on the ultra-widefield set, \rn{indep.DDR.FGADR.drop} on FGADR ($P$ \rn{indep.DDR.FGADR.p}) and \rn{indep.DDR.Retinal-Lesions.drop} on Retinal-Lesions. For models trained on the 44 IDRiD images the losses were larger (\rn{indep.IDRiD.TJDR-std.drop}, \rn{indep.IDRiD.TJDR-uwf.drop}, \rn{indep.IDRiD.FGADR.drop} and \rn{indep.IDRiD.Retinal-Lesions.drop}). The order of the sources was the same for both: least on fine labels at a standard field of view, more on mixed and coarse labels, and most where the imaging geometry changes.

Internal and independent performance were correlated across models (Spearman $\rho$ = \rn{indep.DDR.rho} for DDR), and the model that led on the fine-grained independent tests (\rn{indep.DDR.top.model}, \rn{indep.DDR.top.value}) was \rn{indep.DDR.sametop} as on the internal test. Averaged by family, the self-configuring baseline was ahead on internal (\rn{family.auto.internal}) and on fine-grained independent tests (\rn{family.auto.independent}), fundus-specific networks (\rn{family.fundus.internal} and \rn{family.fundus.independent}) were close to the best general network, and adapted foundation models did not exceed either (Fig.~\ref{fig:retrained}d).

Over the seven test sets the Friedman test rejected equal ranks ($\chi^2$ = \rn{cd.chi2}, $P$ \rn{cd.p}). \rn{cd.first.model} had the best average rank (\rn{cd.first.rank}) and \rn{cd.last.model} the worst (\rn{cd.last.rank}); with seven test sets the Nemenyi critical difference is \rn{cd.value} ranks, and \rn{cd.ntop} of \rn{n.main} models were not separable from the first (Fig.~\ref{fig:retrained}e). The ranking published for the \rn{tau.published.n} models that report IDRiD results agreed only partly with the ranking after retraining (Kendall $\tau$ = \rn{tau.published}, $P$ = \rn{tau.published.p}; Fig.~\ref{fig:retrained}f). Using the source recipes instead of the unified protocol lowered \rn{b1.IDRiD.nlower} of \rn{b1.n} models on IDRiD, by up to \rn{b1.IDRiD.worst.delta} for \rn{b1.IDRiD.worst.model} (Supplementary Table~4), and a threefold change of the learning rate in either direction, tested on DDR for one general and one fundus-specific network, did not change their order (Extended Data Fig.~7).
\end{simtext}

\begin{figure}[p]
\centering
\includegraphics[width=\textwidth,height=0.78\textheight,keepaspectratio]{\figfile{fig4_retrained}}
\caption{\textbf{Performance of retrained models under the unified protocol.} \textbf{a}, Internal tests: per-class AUPR (mean of three seeds) and \mAUPR{} (bar; error bar, s.d.\ over seeds). Models are ordered by their average rank in \textbf{e}; the square gives the family. \textbf{b}, \mAUPR{} of every model on the seven test sets, by training source. Int., internal test; Cross, test set of the other training source; the bar under each column gives the annotation granularity. Columns on the right: internal test, mean of the two fine-grained independent test sets, and their difference. \textbf{c}, The same models on the internal test and on four independent test sets; each line is one model. $P$ values, two-sided Wilcoxon signed-rank tests across models with Holm correction ($n$ = \rn{n.main}). \textbf{d}, \mAUPR{} by family on internal and on fine-grained independent tests; coloured dots, one model and training source; black, mean $\pm$ s.d.; K--W, Kruskal--Wallis test. \textbf{e}, Critical-difference diagram over the seven test sets; bars join models whose average ranks differ by less than the Nemenyi critical difference. \textbf{f}, Rank on IDRiD as published and after retraining.}
\label{fig:retrained}
\end{figure}

\begin{table*}[t]
\centering\tablefont\setlength{\tabcolsep}{3.3pt}
\caption{\textbf{Unified retraining (B-std): independent tests.} The checkpoints of Table~\ref{tab:main} are applied without any adaptation to five test sets from sources never used for training. Values are mAUPR (mean of three seeds); the DiaRetDB1 column is the official image-level ROC AUC, averaged over the four lesion types, because its region-consensus ground truth biases pixel precision. $\Delta$ is the mean of the two fine-grained TJDR test sets minus the internal test. Average rank is over the seven test sets of Fig.~\ref{fig:retrained}e (Friedman $\chi^2$ = \ph{58.1}; Nemenyi critical difference \ph{7.50}).}
\label{tab:independent}
\begin{adjustbox}{max width=\textwidth}
\begin{threeparttable}
\begin{tabular}{@{}l c cccc c c @{\hspace{8pt}} c cccc c c @{\hspace{8pt}} c@{}}
\toprule
 & \multicolumn{7}{c}{\colh{Trained on IDRiD}} & \multicolumn{7}{c}{\colh{Trained on DDR}} & \\
\cmidrule(lr){2-8}\cmidrule(lr){9-15}
 & & \multicolumn{2}{c}{TJDR} & & & & & & \multicolumn{2}{c}{TJDR} & & & & & \colh{Avg.} \\
\colh{Model} & Int. & Std & UWF & FGADR & Ret.-L. & DiaRet & $\Delta$ & Int. & Std & UWF & FGADR & Ret.-L. & DiaRet & $\Delta$ & \colh{rank} \\
 & \tier{fine} & \tier{fine} & \tier{fine} & \tier{mixed} & \tier{coarse} & \tier{region} &  & \tier{fine} & \tier{fine} & \tier{fine} & \tier{mixed} & \tier{coarse} & \tier{region} &  & \\
\midrule
\grouprow{16}{\famchip{general}\; General segmentation}
U-Net & \ph{.66} & \cellcolor[HTML]{D4ECE0}\ph{.23} & \cellcolor[HTML]{EBF4EB}\ph{.10} & \cellcolor[HTML]{E6F3E8}\ph{.13} & \cellcolor[HTML]{E9F3EA}\ph{.11} & \ph{.77} & \cellcolor[HTML]{E47F75}\ph{$$-$$.49} & \ph{.22} & \cellcolor[HTML]{D9EDE2}\ph{.21} & \cellcolor[HTML]{EBF4EB}\ph{.09} & \cellcolor[HTML]{E8F3E9}\ph{.12} & \cellcolor[HTML]{EBF4EB}\ph{.10} & \ph{.77} & \cellcolor[HTML]{EFDFDB}\ph{$$-$$.07} & \ph{12.29} \\
DeepLabv3 & \ph{.54} & \cellcolor[HTML]{CAE8DC}\ph{.27} & \cellcolor[HTML]{E4F2E7}\ph{.16} & \cellcolor[HTML]{E1F1E6}\ph{.18} & \cellcolor[HTML]{E7F3E9}\ph{.13} & \ph{.78} & \cellcolor[HTML]{EAA69E}\ph{$$-$$.32} & \ph{.38} & \cellcolor[HTML]{C3E5D9}\ph{.30} & \cellcolor[HTML]{E2F1E6}\ph{.17} & \cellcolor[HTML]{DCEFE4}\ph{.20} & \cellcolor[HTML]{E5F2E8}\ph{.14} & \ph{.79} & \cellcolor[HTML]{EECFCA}\ph{$$-$$.14} & \ph{8.29} \\
HRNet-W48 & \ph{\best{.70}} & \cellcolor[HTML]{CEE9DD}\ph{.26} & \cellcolor[HTML]{E4F2E7}\ph{.16} & \cellcolor[HTML]{D4ECE0}\ph{.23} & \cellcolor[HTML]{E8F3E9}\ph{.12} & \ph{\second{.86}} & \cellcolor[HTML]{E47F75}\ph{$$-$$.49} & \ph{\best{.45}} & \cellcolor[HTML]{C4E5D9}\ph{.30} & \cellcolor[HTML]{E0F0E5}\ph{.18} & \cellcolor[HTML]{CBE8DC}\ph{.27} & \cellcolor[HTML]{E6F2E8}\ph{.14} & \ph{\second{.87}} & \cellcolor[HTML]{ECC0BA}\ph{$$-$$.21} & \ph{4.43} \\
Swin-Unet & \ph{\second{.69}} & \cellcolor[HTML]{D3EBE0}\ph{.24} & \cellcolor[HTML]{EBF4EB}\ph{.10} & \cellcolor[HTML]{E3F2E7}\ph{.16} & \cellcolor[HTML]{E9F3EA}\ph{.11} & \ph{.76} & \cellcolor[HTML]{E37D74}\ph{$$-$$.52} & \ph{.43} & \cellcolor[HTML]{CBE8DC}\ph{.27} & \cellcolor[HTML]{E9F3EA}\ph{.11} & \cellcolor[HTML]{DFF0E5}\ph{.18} & \cellcolor[HTML]{E8F3E9}\ph{.12} & \ph{.78} & \cellcolor[HTML]{ECB8B1}\ph{$$-$$.24} & \ph{9.43} \\
FCT & \ph{.30} & \cellcolor[HTML]{CBE8DC}\ph{.27} & \cellcolor[HTML]{EAF4EA}\ph{.10} & \cellcolor[HTML]{E5F2E8}\ph{.14} & \cellcolor[HTML]{EAF4EA}\ph{.10} & \ph{.78} & \cellcolor[HTML]{EED5D1}\ph{$$-$$.11} & \ph{.18} & \cellcolor[HTML]{C6E6DA}\ph{.29} & \cellcolor[HTML]{E9F3EA}\ph{.11} & \cellcolor[HTML]{E3F2E7}\ph{.16} & \cellcolor[HTML]{E9F3EA}\ph{.11} & \ph{.79} & \cellcolor[HTML]{EAECEC}\ph{+.03} & \ph{12.14} \\
\grouprow{16}{\famchip{auto}\; Self-configuring}
nnU-Net & \ph{.67} & \cellcolor[HTML]{C3E5D9}\ph{.31} & \cellcolor[HTML]{DFF0E5}\ph{\best{.19}} & \cellcolor[HTML]{D0EADE}\ph{\second{.25}} & \cellcolor[HTML]{D8EDE2}\ph{\best{.22}} & \ph{.85} & \cellcolor[HTML]{E98F85}\ph{$$-$$.42} & \ph{\second{.43}} & \cellcolor[HTML]{BEE3D7}\ph{.33} & \cellcolor[HTML]{DBEEE3}\ph{.20} & \cellcolor[HTML]{CBE8DC}\ph{.27} & \cellcolor[HTML]{D5ECE1}\ph{\best{.22}} & \ph{.86} & \cellcolor[HTML]{EDC8C3}\ph{$$-$$.17} & \ph{\best{2.86}} \\
nnU-Net ResEnc-M & \ph{.63} & \cellcolor[HTML]{AEDDD0}\ph{\best{.39}} & \cellcolor[HTML]{E4F2E7}\ph{.15} & \cellcolor[HTML]{C5E6DA}\ph{\best{.29}} & \cellcolor[HTML]{E3F2E7}\ph{.16} & \ph{\best{.86}} & \cellcolor[HTML]{EA9D95}\ph{$$-$$.36} & \ph{.43} & \cellcolor[HTML]{A1D8CE}\ph{\best{.43}} & \cellcolor[HTML]{E0F0E5}\ph{.18} & \cellcolor[HTML]{BEE3D7}\ph{\best{.33}} & \cellcolor[HTML]{E2F1E6}\ph{.17} & \ph{\best{.88}} & \cellcolor[HTML]{EED4CF}\ph{$$-$$.12} & \ph{\second{3.29}} \\
\grouprow{16}{\famchip{fundus}\; Fundus-specific}
M2MRF & \ph{.66} & \cellcolor[HTML]{C4E5D9}\ph{.30} & \cellcolor[HTML]{E3F1E6}\ph{.16} & \cellcolor[HTML]{D4ECE0}\ph{.23} & \cellcolor[HTML]{E5F2E8}\ph{.14} & \ph{.82} & \cellcolor[HTML]{E98F85}\ph{$$-$$.43} & \ph{.37} & \cellcolor[HTML]{B7E0D4}\ph{.36} & \cellcolor[HTML]{DAEEE3}\ph{\second{.21}} & \cellcolor[HTML]{C8E7DB}\ph{\second{.28}} & \cellcolor[HTML]{E2F1E6}\ph{.17} & \ph{.84} & \cellcolor[HTML]{EFDCD8}\ph{$$-$$.09} & \ph{4.57} \\
HACDR-Net & \ph{.56} & \cellcolor[HTML]{C8E7DB}\ph{.29} & \cellcolor[HTML]{E8F3E9}\ph{.12} & \cellcolor[HTML]{E4F2E7}\ph{.15} & \cellcolor[HTML]{E8F3E9}\ph{.12} & \ph{.79} & \cellcolor[HTML]{EA9D95}\ph{$$-$$.36} & \ph{.27} & \cellcolor[HTML]{BAE2D5}\ph{.34} & \cellcolor[HTML]{E4F2E7}\ph{.15} & \cellcolor[HTML]{DCEFE4}\ph{.20} & \cellcolor[HTML]{E5F2E8}\ph{.14} & \ph{.82} & \cellcolor[HTML]{F0EBE8}\ph{$$-$$.02} & \ph{8.86} \\
\grouprow{16}{\famchip{foundation}\; Foundation models, adapted}
SAMed & \ph{.53} & \cellcolor[HTML]{D3EBE0}\ph{.24} & \cellcolor[HTML]{E7F3E9}\ph{.13} & \cellcolor[HTML]{DCEFE4}\ph{.20} & \cellcolor[HTML]{E5F2E8}\ph{.14} & \ph{.81} & \cellcolor[HTML]{EAA199}\ph{$$-$$.34} & \ph{.29} & \cellcolor[HTML]{CEE9DD}\ph{.26} & \cellcolor[HTML]{E5F2E8}\ph{.15} & \cellcolor[HTML]{D9EDE2}\ph{.21} & \cellcolor[HTML]{E5F2E7}\ph{.15} & \ph{.84} & \cellcolor[HTML]{EFDCD8}\ph{$$-$$.08} & \ph{8.57} \\
SAM-Adapter & \ph{.50} & \cellcolor[HTML]{D8EDE2}\ph{.21} & \cellcolor[HTML]{E7F3E9}\ph{.13} & \cellcolor[HTML]{E1F1E6}\ph{.18} & \cellcolor[HTML]{EAF4EA}\ph{.10} & \ph{.80} & \cellcolor[HTML]{EAA49C}\ph{$$-$$.33} & \ph{.26} & \cellcolor[HTML]{D4ECE0}\ph{.23} & \cellcolor[HTML]{E5F2E8}\ph{.14} & \cellcolor[HTML]{DCEFE4}\ph{.19} & \cellcolor[HTML]{E9F3EA}\ph{.11} & \ph{.80} & \cellcolor[HTML]{EFDEDA}\ph{$$-$$.07} & \ph{11.43} \\
SAM2-UNet & \ph{.61} & \cellcolor[HTML]{BBE2D6}\ph{\second{.34}} & \cellcolor[HTML]{E0F0E5}\ph{\second{.18}} & \cellcolor[HTML]{D9EDE2}\ph{.21} & \cellcolor[HTML]{E4F2E7}\ph{.15} & \ph{.84} & \cellcolor[HTML]{EA9F97}\ph{$$-$$.35} & \ph{.39} & \cellcolor[HTML]{AFDDD1}\ph{\second{.39}} & \cellcolor[HTML]{D8EDE2}\ph{\best{.22}} & \cellcolor[HTML]{CFEADE}\ph{.25} & \cellcolor[HTML]{E0F0E5}\ph{.18} & \ph{.87} & \cellcolor[HTML]{EFDCD8}\ph{$$-$$.08} & \ph{3.86} \\
MedSAM (full FT) & \ph{.51} & \cellcolor[HTML]{CDE9DD}\ph{.26} & \cellcolor[HTML]{E8F3E9}\ph{.12} & \cellcolor[HTML]{D6EDE1}\ph{.22} & \cellcolor[HTML]{E7F3E9}\ph{.13} & \ph{.80} & \cellcolor[HTML]{EAA69E}\ph{$$-$$.32} & \ph{.28} & \cellcolor[HTML]{C8E7DB}\ph{.28} & \cellcolor[HTML]{E6F3E8}\ph{.13} & \cellcolor[HTML]{D0EADE}\ph{.25} & \cellcolor[HTML]{E6F3E8}\ph{.13} & \ph{.81} & \cellcolor[HTML]{EFDFDB}\ph{$$-$$.07} & \ph{9.43} \\
RETFound + UPerNet & \ph{.58} & \cellcolor[HTML]{C2E5D8}\ph{.31} & \cellcolor[HTML]{E3F2E7}\ph{.16} & \cellcolor[HTML]{E0F0E5}\ph{.18} & \cellcolor[HTML]{DCEFE4}\ph{\second{.20}} & \ph{.84} & \cellcolor[HTML]{EAA199}\ph{$$-$$.35} & \ph{.33} & \cellcolor[HTML]{BDE3D6}\ph{.33} & \cellcolor[HTML]{E1F1E6}\ph{.18} & \cellcolor[HTML]{DBEEE3}\ph{.20} & \cellcolor[HTML]{DAEEE3}\ph{\second{.21}} & \ph{.86} & \cellcolor[HTML]{EFDCD8}\ph{$$-$$.08} & \ph{5.57} \\
\bottomrule
\end{tabular}
\begin{tablenotes}[flushleft]\scriptsize
\item Granularity of the test set: \tier{fine} fine \quad \tier{mixed} mixed \quad \tier{coarse} coarse \quad \tier{region} region consensus. Int., internal test (official test set of the training source). Std, standard field (35--50$^\circ$); UWF, ultra-widefield (133$^\circ$), predicted at native resolution with sliding windows. Retinal-Lesions excludes its grey-level 127 regions.
\phnote
\end{tablenotes}
\end{threeparttable}
\end{adjustbox}
\end{table*}


% ============================================================== R5
\subsection{Foundation models under four kinds of use}
Results for segmentation foundation models are often quoted without saying how the model was used. We separated four kinds of use and report them side by side without ranking one against another (Fig.~\ref{fig:foundation}a and Table~\ref{tab:foundation}): zero-shot use without any label; prompts derived from the ground truth, which is an upper bound and not a usable setting; parameter-efficient fine-tuning; and an encoder with a trained decoder or full fine-tuning.

\begin{simtext}
Zero-shot use did not segment DR lesions. Automatic masks of SAM and SAM~2, even when each mask was given the class it overlaps most, and text prompts of SAM~3 stayed far below every trained model (best: \rn{fm.zero.model}, \mAUPR{} \rn{fm.zero.value} on IDRiD). Text prompts were sensitive to wording: across three phrasings per lesion the best and worst AUPR differed \rn{fm.wording.ratio}-fold on average (Fig.~\ref{fig:foundation}c). With a box around every ground-truth component SAM~2 reached \rn{fm.oracle.all}, but an IDRiD image carries on average 44 microaneurysm and 142 hard exudate components; with 5 and 20 boxes per image and class the \mAUPR{} was \rn{fm.oracle.five} and \rn{fm.oracle.twenty} (Fig.~\ref{fig:foundation}b). Given the same boxes, the supervised models drew the lesions at least as well inside the box, so the gap of the promptable models lies in finding lesions, not in delineating them.

Among the trained settings, \rn{fm.peft.model} was the best parameter-efficient adaptation (\rn{fm.peft.value} on DDR, against \rn{fm.ref.nnunet} for nnU-Net). Under one UPerNet head, one recipe and one training amount on DDR, an encoder pretrained on retinal images (\rn{fm.enc.retfound.int} internal, \rn{fm.enc.retfound.ind} independent) was compared with its own starting point, a masked-autoencoder trained on ImageNet (\rn{fm.enc.imagenet.int} and \rn{fm.enc.imagenet.ind}), and the same contrast was repeated under the DINOv2 objective (retinal \rn{fm.enc.retfound.ind} against general \rn{fm.enc.dinov2.ind} on independent data), so that the effect of retinal pretraining does not depend on one pretraining objective (Fig.~\ref{fig:foundation}d,e).
\end{simtext}

\begin{figure}[t]
\centering
\widefig{fig5_foundation}
\caption{\textbf{Foundation models under four kinds of use.} \textbf{a}, Per-class AUPR and \mAUPR{} on the IDRiD and DDR test sets. Zero-shot and oracle rows use no training image; the two trained settings are trained on the source of the test set. \textbf{b}, Oracle prompts: AUPR on IDRiD against the number of ground-truth boxes or points given per image and class; dotted line, nnU-Net trained on IDRiD. \textbf{c}, SAM~3 with three wordings per lesion class. \textbf{d}, Four ViT-L encoders under the same UPerNet head and recipe, trained on DDR (seed 0): internal test and mean of the two TJDR test sets; lines join each retinal encoder to the natural-image encoder trained with the same objective. \textbf{e}, Best model of each kind of use against the fundus-specific and self-configuring references.}
\label{fig:foundation}
\end{figure}

\begin{table*}[t]
\centering\tablefont
\caption{\textbf{Foundation models under four kinds of use.} The kinds are reported in separate blocks and are not ranked against each other: zero-shot use sees no test label (its density, area cap and wording are chosen on the validation pool); the upper-bound block reads the test labels (masks classed by their overlap with the ground truth, prompts derived from every ground-truth component) and is not a usable setting; the two trained settings follow the unified protocol. AUPR on the IDRiD and DDR test sets. \best{Bold} and \second{underlined} mark the best and second-best rows of each column within a block.}
\label{tab:foundation}
\begin{adjustbox}{max width=\textwidth}
\begin{threeparttable}
\begin{tabular}{@{}l l l cccc c @{\hspace{10pt}} cccc c@{}}
\toprule
 & & & \multicolumn{5}{c}{\colh{IDRiD test}} & \multicolumn{5}{c}{\colh{DDR test}} \\
\cmidrule(lr){4-8}\cmidrule(l){9-13}
\colh{Model} & \colh{Use} & \colh{Backbone} & MA & HE & EX & SE & \colh{mAUPR} & MA & HE & EX & SE & \colh{mAUPR} \\
\midrule
\grouprow{13}{Zero-shot: no test label; density, area cap and wording chosen on the validation pool (summary = any-lesion AUPR$^\ddagger$)}
\famchip{foundation}\; SAM & automatic masks, class-agnostic & ViT-H & \na & \na & \na & \na & \ph{.06} & \na & \na & \na & \na & \ph{.06} \\
\famchip{foundation}\; SAM 2 & automatic masks, class-agnostic & Hiera-L & \na & \na & \na & \na & \ph{\second{.07}} & \na & \na & \na & \na & \ph{\second{.07}} \\
\famchip{foundation}\; SAM 3 & text, wording chosen on validation & -- & \cellcolor[HTML]{F3F7EF}\ph{.03} & \cellcolor[HTML]{E5F2E7}\ph{.19} & \cellcolor[HTML]{D6EDE1}\ph{.29} & \cellcolor[HTML]{E2F1E6}\ph{.22} & \ph{\best{.16}} & \cellcolor[HTML]{F3F7EF}\ph{.03} & \cellcolor[HTML]{E9F4EA}\ph{.14} & \cellcolor[HTML]{DEEFE4}\ph{.25} & \cellcolor[HTML]{E3F1E6}\ph{.21} & \ph{\best{.12}} \\
\grouprow{13}{Upper bounds that read the test labels (not a usable setting)}
\famchip{foundation}\; SAM & automatic masks, class by overlap$^\dagger$ & ViT-H & \cellcolor[HTML]{F3F7F0}\ph{.02} & \cellcolor[HTML]{EBF4EB}\ph{.12} & \cellcolor[HTML]{E6F2E8}\ph{.18} & \cellcolor[HTML]{EAF4EA}\ph{.13} & \ph{.11} & \cellcolor[HTML]{F3F7F0}\ph{.02} & \cellcolor[HTML]{EBF4EB}\ph{.12} & \cellcolor[HTML]{E8F3E9}\ph{.15} & \cellcolor[HTML]{E9F4EA}\ph{.13} & \ph{.11} \\
\famchip{foundation}\; SAM 2 & automatic masks, class by overlap$^\dagger$ & Hiera-L & \cellcolor[HTML]{F3F7F0}\ph{.02} & \cellcolor[HTML]{E8F3E9}\ph{.15} & \cellcolor[HTML]{E8F3E9}\ph{.16} & \cellcolor[HTML]{E0F0E5}\ph{.23} & \ph{.14} & \cellcolor[HTML]{F2F7EF}\ph{.04} & \cellcolor[HTML]{E9F4EA}\ph{.14} & \cellcolor[HTML]{E2F1E6}\ph{.22} & \cellcolor[HTML]{DEEFE4}\ph{.25} & \ph{.16} \\
\famchip{foundation}\; SAM & oracle boxes, every component & ViT-B & \cellcolor[HTML]{D4ECE0}\ph{\second{.30}} & \cellcolor[HTML]{93D2CB}\ph{\second{.61}} & \cellcolor[HTML]{87CDC8}\ph{\second{.66}} & \cellcolor[HTML]{77C7C5}\ph{\second{.72}} & \ph{\second{.57}} & \cellcolor[HTML]{DBEEE3}\ph{\second{.26}} & \cellcolor[HTML]{A9DBCF}\ph{\second{.53}} & \cellcolor[HTML]{9AD5CC}\ph{\second{.58}} & \cellcolor[HTML]{8CD0C9}\ph{\second{.64}} & \ph{\second{.50}} \\
\famchip{foundation}\; SAM 2 & oracle points, every component & Hiera-L & \cellcolor[HTML]{E3F1E6}\ph{.21} & \cellcolor[HTML]{AFDDD1}\ph{.50} & \cellcolor[HTML]{ACDCD0}\ph{.52} & \cellcolor[HTML]{91D2CA}\ph{.62} & \ph{.46} & \cellcolor[HTML]{E5F2E8}\ph{.19} & \cellcolor[HTML]{B9E1D5}\ph{.45} & \cellcolor[HTML]{B8E1D4}\ph{.45} & \cellcolor[HTML]{A0D7CD}\ph{.56} & \ph{.41} \\
\famchip{foundation}\; SAM 2 & oracle boxes, every component & Hiera-L & \cellcolor[HTML]{CDE9DD}\ph{\best{.34}} & \cellcolor[HTML]{88CEC9}\ph{\best{.65}} & \cellcolor[HTML]{7CC9C6}\ph{\best{.70}} & \cellcolor[HTML]{6EC4C3}\ph{\best{.75}} & \ph{\best{.61}} & \cellcolor[HTML]{D6EDE1}\ph{\best{.29}} & \cellcolor[HTML]{99D5CC}\ph{\best{.59}} & \cellcolor[HTML]{93D2CB}\ph{\best{.61}} & \cellcolor[HTML]{85CDC8}\ph{\best{.67}} & \ph{\best{.54}} \\
\famchip{foundation}\; MedSAM & oracle boxes, every component & ViT-B & \cellcolor[HTML]{E3F2E7}\ph{.20} & \cellcolor[HTML]{9ED7CD}\ph{.57} & \cellcolor[HTML]{9CD6CC}\ph{.58} & \cellcolor[HTML]{7ECAC6}\ph{.69} & \ph{.51} & \cellcolor[HTML]{E4F2E7}\ph{.20} & \cellcolor[HTML]{AFDDD1}\ph{.51} & \cellcolor[HTML]{AEDDD0}\ph{.51} & \cellcolor[HTML]{90D1CA}\ph{.62} & \ph{.46} \\
\grouprow{13}{Parameter-efficient fine-tuning (trained on the source of the test set)}
\famchip{foundation}\; SAMed & LoRA & ViT-B & \cellcolor[HTML]{E1F1E6}\ph{\second{.23}} & \cellcolor[HTML]{A9DBCF}\ph{\second{.53}} & \cellcolor[HTML]{6EC4C3}\ph{\second{.75}} & \cellcolor[HTML]{A7DACF}\ph{\second{.54}} & \ph{\second{.51}} & \cellcolor[HTML]{ECF5EC}\ph{\second{.11}} & \cellcolor[HTML]{C8E7DB}\ph{\second{.37}} & \cellcolor[HTML]{B8E1D4}\ph{\second{.46}} & \cellcolor[HTML]{D8EDE2}\ph{\second{.28}} & \ph{\second{.30}} \\
\famchip{foundation}\; SAM-Adapter & adapters & ViT-H & \cellcolor[HTML]{E3F1E6}\ph{.21} & \cellcolor[HTML]{AEDDD0}\ph{.51} & \cellcolor[HTML]{73C6C4}\ph{.73} & \cellcolor[HTML]{AADCCF}\ph{.52} & \ph{.49} & \cellcolor[HTML]{EDF5EC}\ph{.09} & \cellcolor[HTML]{CFEADE}\ph{.33} & \cellcolor[HTML]{BEE3D7}\ph{.42} & \cellcolor[HTML]{DCEFE4}\ph{.25} & \ph{.27} \\
\famchip{foundation}\; SAM2-UNet & adapters + U-shaped decoder & Hiera-L & \cellcolor[HTML]{C6E6DA}\ph{\best{.38}} & \cellcolor[HTML]{8ACFC9}\ph{\best{.65}} & \cellcolor[HTML]{5DBAC0}\ph{\best{.82}} & \cellcolor[HTML]{88CEC9}\ph{\best{.65}} & \ph{\best{.62}} & \cellcolor[HTML]{E5F2E8}\ph{\best{.18}} & \cellcolor[HTML]{BAE2D5}\ph{\best{.44}} & \cellcolor[HTML]{A3D9CE}\ph{\best{.55}} & \cellcolor[HTML]{CAE8DC}\ph{\best{.35}} & \ph{\best{.38}} \\
\grouprow{13}{Encoder + decoder (trained on the source of the test set)}
\famchip{foundation}\; MedSAM & full fine-tuning & ViT-B & \cellcolor[HTML]{E2F1E6}\ph{\second{.22}} & \cellcolor[HTML]{A7DACF}\ph{\second{.54}} & \cellcolor[HTML]{6CC3C3}\ph{\second{.76}} & \cellcolor[HTML]{A1D8CE}\ph{\second{.56}} & \ph{\second{.52}} & \cellcolor[HTML]{EDF5EC}\ph{\second{.10}} & \cellcolor[HTML]{CBE8DC}\ph{\second{.35}} & \cellcolor[HTML]{B8E1D4}\ph{\second{.46}} & \cellcolor[HTML]{D8EDE2}\ph{\second{.28}} & \ph{\second{.30}} \\
\famchip{foundation}\; RETFound & MAE encoder + UPerNet & ViT-L & \cellcolor[HTML]{D8EDE2}\ph{\best{.28}} & \cellcolor[HTML]{99D5CC}\ph{\best{.59}} & \cellcolor[HTML]{68C1C2}\ph{\best{.77}} & \cellcolor[HTML]{97D4CB}\ph{\best{.60}} & \ph{\best{.56}} & \cellcolor[HTML]{EBF4EB}\ph{\best{.12}} & \cellcolor[HTML]{C6E6DA}\ph{\best{.38}} & \cellcolor[HTML]{B3DFD2}\ph{\best{.48}} & \cellcolor[HTML]{D4ECE0}\ph{\best{.30}} & \ph{\best{.32}} \\
\midrule
\famchip{fundus}\; M2MRF & reference, Table~\ref{tab:main} & -- & \ph{.45} & \ph{.63} & \ph{.87} & \ph{.68} & \ph{.66} & \ph{.15} & \ph{.41} & \ph{.53} & \ph{.39} & \ph{.37} \\
\famchip{auto}\; nnU-Net & reference, Table~\ref{tab:main} & -- & \ph{.50} & \ph{.67} & \ph{.87} & \ph{.62} & \ph{.66} & \ph{.16} & \ph{.50} & \ph{.61} & \ph{.45} & \ph{.43} \\
\bottomrule
\end{tabular}
\begin{tablenotes}[flushleft]\scriptsize
\item $^\ddagger$The zero-shot block is summarised by the any-lesion AUPR (union of the four classes against the pixelwise maximum of the maps); automatic masks carry no class, so their per-class cells are empty, and SAM~3's per-class cells are its own text prompts. The any-lesion AUPR of the references on the same test sets (IDRiD / DDR): M2MRF -- / --; nnU-Net 0.78 / 0.57 (-- where the run predates the metric).
\item $^\dagger$Each automatic mask is assigned the class it overlaps most in the ground truth, an upper bound for class-agnostic proposals. Scores of automatic masks use the predicted IoU.
\phnote
\end{tablenotes}
\end{threeparttable}
\end{adjustbox}
\end{table*}


% ============================================================== R6
\subsection{Effects of annotation granularity and training source}
The sources differ in how lesions were drawn before they differ in anything a model does (Fig.~\ref{fig:sources}a). In the fine-grained sources an image carries many small components; in the coarse and region-consensus sources it carries few and large ones, which for the same retina means that several lesions and the tissue between them were enclosed by one outline.

\begin{simtext}
To separate the effect of the training source from the effect of the model, the five representative models were also trained on the standard-field TJDR images and on FGADR, at the pixel budget of DDR and with one seed, which together with the main tables gives a matrix of four training sources by seven test sets (Fig.~\ref{fig:sources}b,d). Every test set was served best by its own source. Training elsewhere cost \rn{cross.mean} \mAUPR{} on average and up to \rn{cross.max} (Fig.~\ref{fig:sources}c); models trained on the mixed labels of FGADR lost \rn{cross.fgadr.fine} on the three fine-grained sources.

We then added mixed or coarse labels to fine-label training while keeping the number of image passes fixed, for U-Net and M2MRF with one seed; differences smaller than twice the median seed standard deviation of the main table are reported without interpretation (Fig.~\ref{fig:sources}e and Table~\ref{tab:factors}). Relative to DDR alone (\rn{e4.ref} on the IDRiD test set), adding FGADR changed \mAUPR{} by \rn{e4.fgadr}; adding only 383 FGADR images, as many as DDR has, by \rn{e4.fgadr383}, which separates the effect of more images from the effect of a different label style. Adding Retinal-Lesions changed it by \rn{e4.retles}, and training on Retinal-Lesions or FGADR alone by \rn{e4.retlesonly} and \rn{e4.fgadronly}.
\end{simtext}

\begin{figure}[t]
\centering
\widefig{fig6_sources}
\caption{\textbf{Effects of annotation granularity and training source.} \textbf{a}, Annotation style: lesions per image against the median lesion diameter relative to the field-of-view diameter, per lesion class and source. \textbf{b}, \mAUPR{} for each training source (rows) and test set (columns), mean of five representative models, seed 0; outlined cells are internal tests. \textbf{c}, Difference to training on the source of the test set. \textbf{d}, The matrix of \textbf{b} for each model. \textbf{e}, Training sets of different granularity at an equal number of image passes, tested on two fine-grained test sets (U-Net and M2MRF, seed 0).}
\label{fig:sources}
\end{figure}

% ============================================================== R7
\subsection{Effects of lesion size, input resolution and training amount}
Lesion size is the property that separates the four classes most clearly (Fig.~\ref{fig:scale}a). On the training images the median diameter of a microaneurysm is 0.5\% of the field-of-view diameter and that of a hemorrhage 1.1\%; we froze these two values, computed on training images only, as the boundaries between small, medium and large lesions before the main-table training. Resizing a whole image to a small input removes lesions from the training target itself: at 224 pixels \rn{vanish.IDRiD.224.MA}\% of the microaneurysms and \rn{vanish.IDRiD.224.EX}\% of the hard exudates of the IDRiD training set no longer cover a single pixel, and on DDR \rn{vanish.DDR.224.MA}\% and \rn{vanish.DDR.224.EX}\%; at 384 pixels the DDR figures are \rn{vanish.DDR.384.MA}\% and \rn{vanish.DDR.384.EX}\%, and on the 1,440-pixel canvas of the unified protocol essentially none are lost (Fig.~\ref{fig:scale}b).

\begin{simtext}
Lesion-level recall followed size for every model (Fig.~\ref{fig:scale}c). On DDR the median recall over models was \rn{recall.DDR.small.median} for small and \rn{recall.DDR.large.median} for large lesions, and the best model detected \rn{recall.DDR.small.best} of the small ones. Input resolution acted mainly through this channel. For M2MRF and for the RETFound encoder with a UPerNet head, \mAUPR{} on DDR rose from \rn{e7.aupr.512} at a long side of 512 pixels to \rn{e7.aupr.1440} at 1,440 and \rn{e7.aupr.native} with native-resolution patches (mean of the two models, one seed), while the recall of small lesions rose from \rn{e7.small.512} to \rn{e7.small.1440} (Fig.~\ref{fig:scale}d).

More training images helped up to the size of the available sources. With 10\%, 50\% and 100\% of the DDR training set, \mAUPR{} on the DDR test set was \rn{e8.ddr.10}, \rn{e8.ddr.50} and \rn{e8.ddr.100} (U-Net and M2MRF, one seed, equal number of image passes), and merging the three fine-grained sources IDRiD, DDR and standard-field TJDR gave \rn{e8.ddr.all}. A source added to the training set raised the score on its own test set (TJDR standard field: \rn{e8.tjdr.100} before and \rn{e8.tjdr.fine} after it was added), which is the gain of in-source data seen in the matrix above and not a gain in generalization (Fig.~\ref{fig:scale}e). Exchanging the loss of U-Net with the architecture fixed moved \mAUPR{} by at most \rn{e9.range}, except for the per-image Dice term of the registered recipe, which is the collapse described above (Extended Data Fig.~8).
\end{simtext}

\begin{figure}[p]
\centering
\widefig{fig7_scale}
\caption{\textbf{Effects of lesion size, input resolution and training amount.} \textbf{a}, Lesion diameter relative to the field-of-view diameter on the training images; dashed lines, the frozen boundaries between small, medium and large lesions. \textbf{b}, Share of connected components that vanish from the training target when the whole image is resized to the given size. \textbf{c}, Lesion-level recall by lesion size and lesion-level precision on the internal tests, at the threshold chosen on the validation set. \textbf{d}, \mAUPR{} and recall of small lesions on DDR against the long side of the canvas (M2MRF and RETFound with UPerNet, seed 0). \textbf{e}, \mAUPR{} against the fraction of the DDR training set and after the fine-grained sources are merged in (U-Net and M2MRF, seed 0); the number of image passes is the same in every setting.}
\label{fig:scale}
\end{figure}

\begin{table*}[t]
\centering\tablefont
\caption{\textbf{One factor at a time.} One or two representative models are retrained with a single factor changed and everything else fixed to the unified protocol (seed 0). Each cell is the value followed by the change against the reference row of its block; -- marks a model the block does not use. The last column averages the models of the block; its shading follows the change. Differences smaller than twice the median seed standard deviation of the main table are not interpreted.}
\label{tab:factors}
\begin{adjustbox}{max width=\textwidth}
\begin{threeparttable}
\begin{tabular}{@{}l ccc c@{}}
\toprule
\colh{Setting} & \famchip{general}\; U-Net & \famchip{fundus}\; M2MRF & \famchip{foundation}\; RETFound+UPer. & \colh{Mean} \\
\midrule
\grouprow{5}{Annotation granularity of the training set (E4) --- mAUPR on IDRiD test, equal image passes; U-Net and M2MRF}
\quad DDR \tier{fine} {\scriptsize\color{inkmuted}(ref.)} & \ph{.36} & \ph{.54} & -- & \ph{.45} \\
\quad DDR + FGADR, 383 images \tier{fine}\tier{mixed} & \ph{.36} \same{\ph{.00}} & \ph{.54} \up{\ph{.01}} & -- & \cellcolor[HTML]{EFEFEC}\ph{.45} {\scriptsize \ph{+.00}} \\
\quad DDR + FGADR \tier{fine}\tier{mixed} & \ph{.38} \up{\ph{.02}} & \ph{.59} \up{\ph{.06}} & -- & \cellcolor[HTML]{E2E7EC}\ph{.48} {\scriptsize \ph{+.04}} \\
\quad DDR + Retinal-Lesions \tier{fine}\tier{coarse} & \ph{.32} \down{\ph{.04}} & \ph{.52} \down{\ph{.02}} & -- & \cellcolor[HTML]{EFE4E1}\ph{.42} {\scriptsize \ph{$$-$$.03}} \\
\quad FGADR \tier{mixed} & \ph{.27} \down{\ph{.09}} & \ph{.46} \down{\ph{.08}} & -- & \cellcolor[HTML]{EECFCA}\ph{.36} {\scriptsize \ph{$$-$$.09}} \\
\quad Retinal-Lesions \tier{coarse} & \ph{.21} \down{\ph{.15}} & \ph{.35} \down{\ph{.18}} & -- & \cellcolor[HTML]{EBB0A8}\ph{.28} {\scriptsize \ph{$$-$$.17}} \\
\grouprow{5}{Input resolution, long side in pixels (E7) --- mAUPR on DDR test; M2MRF and RETFound + UPerNet}
\quad 512 & -- & \ph{.23} \down{\ph{.19}} & \ph{.23} \down{\ph{.12}} & \cellcolor[HTML]{EBB3AC}\ph{.23} {\scriptsize \ph{$$-$$.16}} \\
\quad 1,440 {\scriptsize\color{inkmuted}(ref.)} & -- & \ph{.42} & \ph{.35} & \ph{.39} \\
\quad native-resolution patches & -- & \ph{.41} \down{\ph{.01}} & \ph{.34} \down{\ph{.01}} & \cellcolor[HTML]{F0EBE8}\ph{.38} {\scriptsize \ph{$$-$$.01}} \\
\grouprow{5}{Input resolution (E7) --- lesion-level recall of small lesions on DDR test}
\quad 512 & -- & \ph{.05} \down{\ph{.29}} & \ph{.11} \down{\ph{.20}} & \cellcolor[HTML]{E99289}\ph{.08} {\scriptsize \ph{$$-$$.24}} \\
\quad 1,440 {\scriptsize\color{inkmuted}(ref.)} & -- & \ph{.34} & \ph{.31} & \ph{.32} \\
\quad native-resolution patches & -- & \ph{.39} \up{\ph{.05}} & \ph{.29} \down{\ph{.01}} & \cellcolor[HTML]{E9EBEC}\ph{.34} {\scriptsize \ph{+.02}} \\
\grouprow{5}{Training amount (E8) --- mAUPR on DDR test; U-Net and M2MRF}
\quad DDR, 10\% (38 images) & \ph{.18} \down{\ph{.16}} & \ph{.23} \down{\ph{.17}} & -- & \cellcolor[HTML]{EBB1AA}\ph{.20} {\scriptsize \ph{$$-$$.16}} \\
\quad DDR, 50\% (192) & \ph{.29} \down{\ph{.05}} & \ph{.35} \down{\ph{.04}} & -- & \cellcolor[HTML]{EFDEDA}\ph{.32} {\scriptsize \ph{$$-$$.05}} \\
\quad DDR, 100\% (383) {\scriptsize\color{inkmuted}(ref.)} & \ph{.33} & \ph{.39} & -- & \ph{.36} \\
\quad + IDRiD + TJDR standard field (599) & \ph{.37} \up{\ph{.04}} & \ph{.44} \up{\ph{.04}} & -- & \cellcolor[HTML]{E0E7EC}\ph{.40} {\scriptsize \ph{+.04}} \\
\grouprow{5}{Loss, architecture fixed (E9) --- mAUPR on DDR test; U-Net}
\quad cross-entropy {\scriptsize\color{inkmuted}(ref.)} & \ph{.30} & -- & -- & \ph{.30} \\
\quad batch Dice & \ph{.34} \up{\ph{.04}} & -- & -- & \cellcolor[HTML]{E0E7EC}\ph{.34} {\scriptsize \ph{+.04}} \\
\quad cross-entropy + batch Dice (unified protocol) & \ph{.34} \up{\ph{.04}} & -- & -- & \cellcolor[HTML]{DFE6EC}\ph{.34} {\scriptsize \ph{+.04}} \\
\quad cross-entropy + per-image Dice (registered recipe) & \ph{.12} \down{\ph{.18}} & -- & -- & \cellcolor[HTML]{EBABA3}\ph{.12} {\scriptsize \ph{$$-$$.18}} \\
\bottomrule
\end{tabular}
\begin{tablenotes}[flushleft]\scriptsize
\item Granularity of the labels used for training: \tier{fine} fine \quad \tier{mixed} mixed \quad \tier{coarse} coarse. In the training-amount block the number of image passes is the same in every row. The FGADR row of the granularity block is the FGADR-trained model of the cross-source experiment.
\phnote
\end{tablenotes}
\end{threeparttable}
\end{adjustbox}
\end{table*}


% ============================================================== R8
\subsection{Computational cost}
\begin{simtext}
All models were trained on single 40\,GB cards, one task per card. Because the training amount is fixed in pixels, the cost of a model under the unified protocol reflects its throughput and its patch size: one seed on IDRiD took from \rn{cost.min} GPU hours (\rn{cost.min.model}) to \rn{cost.max} (\rn{cost.max.model}), and a DDR seed about four times as long, the slowest being the models that process small patches with sliding windows (FCT, 72 hours on DDR). Accuracy was not related to cost or to parameter count (Extended Data Fig.~10 and Supplementary Table~7).
\end{simtext}

\FloatBarrier

\section{Discussion}

We assessed methods for segmenting diabetic retinopathy lesions in fundus photographs by auditing published numbers, scoring and rerunning the original models, retraining \rn{n.main} models under one protocol, and testing all of them on sources they had never seen. The design separates three things that a single leaderboard number mixes: what a paper reports, what its code produces, and what its architecture contributes once data, budget and evaluation are equal.

\paragraph{Reproducibility (Q1).} Most fundus-specific methods cannot be rerun: \rn{funnel.train} of \rn{funnel.dr} have training code. Where code and weights exist, the authors' evaluation code reproduced the released numbers exactly, so the distance between a paper and our protocol is a property of the metric and not of the model; but the released training recipe, run with the authors' settings on three seeds, produced models 5--7 points below the released weights. The literature audit shows why published tables cannot be read as rankings: most papers do not say how the AUPR was computed, the fixed-threshold implementation adds a model-dependent amount, and nearly half of the baseline numbers (\rn{lit.base.copied1}\% identical to the method's own paper, \rn{lit.base.copied2}\% to an earlier copy) were carried from one paper to the next, at times with transposed columns.

\paragraph{Comparison (Q2).} Under equal conditions two adapted foundation models (a low-rank adaptation of SAM and RETFound with a segmentation head) and the general-purpose HRNet-W48 were at the top of both internal tests, SAM-Adapter, Swin-Unet and the self-configuring nnU-Net followed within 0.04, and the seven fundus-specific networks were below them, in line with what has been found for three-dimensional medical segmentation\cite{isensee2024revisited}; over the eight test sets \rn{cd.ntop} of the \rn{n.ranked} ranked models could not be separated from the first. That statement holds for the budget and the 1,440-pixel canvas of the unified protocol: the strongest fundus-specific method, HRDecoder, matched the leading models when run with its authors' released weights, which were trained at twice the resolution and on several times the pixels, and a lesion-specific decoder may need exactly the conditions a shared protocol takes away. The differences among the leading models were of the size of the seed-to-seed variation, whereas the loss of a single rare class, or a recipe that does not converge within the budget, moved a model by ten times as much. Outside the training source the order was less stable than the internal tables suggest: internal rank predicted independent rank on the sources whose lesions are drawn the same way, and the lead on coarsely labelled, region-labelled or ultra-widefield images went to models from the lower half of the internal tables. A model chosen on IDRiD or DDR alone is therefore chosen for the labelling style of that source as much as for the lesions. Where one score is needed, averaging the probability maps of several retrained models gained 0.02--0.06 \mAUPR{} over the best single model on every fine-grained test set and $-$0.01 to 0.02 on the coarse and region-labelled ones, more than the differences among the leading architectures; the four adapted foundation models averaged by themselves gave the same gain. On the one source with several independent expert annotations, the agreement of the better models with the expert consensus was within the agreement among the experts, with the caveats that a region consensus is a coarse ground truth and that the experts' confidence maps give only a four-point curve.
The strongest single models of the study were not in the ranked table. Under one decoder head, one recipe and one training amount, RETFound's two encoders scored no higher than the natural-image encoders they were initialised from ($-$\rn{fm.enc.c1.int.abs} and $-$\rn{fm.enc.c2.int.abs} \mAUPR{} on DDR, $-$\rn{fm.enc.i.c1.int.abs} and $-$\rn{fm.enc.i.c2.int.abs} on IDRiD, one seed each), and the ImageNet masked autoencoder and DINOv2, ViT-L encoders that have never seen a fundus photograph, reached \rn{fm.enc.mi.int} and \rn{fm.enc.di.int} on DDR and \rn{fm.enc.i.mi.int} and \rn{fm.enc.i.di.int} on IDRiD under that head, above every ranked model on both sources. These cells were trained as controls after the ranked rows had been registered, each with one seed and one learning rate, so they do not enter the ranks; but they say that the lead of a ViT-L encoder with a UPerNet head over the convolutional and fundus-specific networks is a lead of the architecture and of the scale of pretraining behind the encoder, which retinal pretraining did not add to, and that the leaderboard of this task would be headed by a generic encoder and a generic decoder if such a configuration were registered. Confirming this with more seeds and a learning-rate sweep is the first item of future work. Zero-shot use of promptable models is not a baseline for this task, and results obtained with prompts from the ground truth should be labelled as upper bounds.

\paragraph{What decides performance (Q3).} The largest effects came from the data and the protocol rather than from the architecture. Whether the test images share the annotation style and the field of view of the training images moved \mAUPR{} by 0.16 on average and up to 0.28 between the four training sources, as much as the differences between the architectures of the main table: models trained on DDR kept or improved their score on the other pixel-labelled standard-field sources and lost a third to a half of it where the labels are coarse or regional or the field is ultra-wide, and models trained on the 44 IDRiD images lost on every other source, as Playout and Cheriet observed for U-Net-type baselines\cite{playout2024crossdatasetgeneralization}. Mixed or coarse labels added to fine ones did not help and often hurt, and the most useful foreign source was the largest fine-grained one rather than the largest one. Shrinking the canvas from 1,440 to 512 pixels cost up to 0.07 \mAUPR{} and between a third and three fifths of the small lesions, more images of the same source helped with diminishing returns, and merging the fine-grained sources improved the scores outside them. An internal test therefore says little about new images unless they are labelled like the training set, a conclusion shared with assessments in other fields\cite{lu2026tcrepitopenaturemethods}.

\paragraph{Recommendations.} For developers of methods: (1) use the official split, and publish the list of image identifiers when a new split is needed; (2) report the exact pooled AUPR at the original resolution inside the field of view, and name the implementation; (3) rerun every baseline or state for each number where it was taken from; (4) include at least one test set from a source that was not used for training, and do not average across annotation granularities; (5) include a self-configuring baseline and report several seeds; (6) release probability maps, so that any metric can be recomputed. For users: the lesion-level recall for small lesions and the result on an independent source with the user's field of view are more informative than a mean AUPR on IDRiD.

\paragraph{Limitations and scope.} The scope of the study was fixed before the factor experiments were run, to what one protocol and one cluster could carry: the main table has three seeds for nine of its \rn{n.main} configurations and one seed for the others, the factor experiments (training source, granularity, resolution, training amount, loss) use one to five representative models and one seed and report only differences larger than the seed variation of the main table, the encoder comparison has one seed and one learning rate per cell, and the source-recipe track was stopped before any run finished. Bootstrap intervals over test images are given per class (Supplementary Table~4); an interval for the class mean, from the same replicates, is left for the final version, and on both internal tests the per-class intervals of the leading models overlap. Left for future work are the remaining seeds and models of each factor experiment, a registered ranked row for the natural-image encoders under the UPerNet head and the full encoder grid (supervised ImageNet, frozen encoders, learning-rate sweeps, the adaptation-by-initialization square), the second-tier and rewritten models (UNet++, UPerNet, PerSAM, H2Former, L-Seg, SAA, RTNet), the released WSRFNet weights, vessel maps and text as additional inputs, the intermediate settings of each factor (a 1,024-pixel canvas, a quarter of DDR, merging all sources, the bin loss), and the test sets that were not available at the time (e-ophtha, ROC, ultra-widefield OCTA).  Region-consensus and coarse ground truth bias pixel metrics, which is why DiaRetDB1 is reported with its image-level protocol and Retinal-Lesions as a separate column. The licence of FGADR does not allow derived data to be released, and we do not release them for Retinal-Lesions either. Most sources carry no patient identifiers, so leakage can be audited only through near-duplicate images. The ultra-widefield test set comes from a single centre and device. IDRiD has 27 test images, and intervals on it are wide. The first round used the hyperparameters of each source recipe, with a small learning-rate sweep on DDR as the only tuning. H2Former, although registered in tier~1, did not fit the cards used here at its native input size and is not in the main tables. Finally, the frozen lesion-matching rule is permissive (any overlap counts), so lesion-level precision is reported together with a stricter variant in the supplementary tables.

\section{Methods}

\subsection{Workflow of model evaluation}
The study has three parts that correspond to the three questions. (Q1) Published numbers are audited (the paper audit and the reproducibility funnel), released checkpoints are scored on official and independent test sets, and the authors' pipelines are rerun (author-recipe reruns). (Q2) Models with complete training code are retrained under one protocol on IDRiD and on DDR (unified retraining), including nnU-Net through its own entry point and foundation models under four kinds of use; all checkpoints are applied to six independent test sets (independent tests). (Q3) One factor is varied at a time on representative models (U-Net, nnU-Net, M2MRF, SAM2-UNet, RETFound with a UPerNet head): training source, annotation granularity, input resolution, training amount and loss; the HRDecoder case study of the main text applies the same machinery to one model's own geometry. Computational cost is recorded for every run. Each experiment is named in the text by what it does; the identifiers under which the runs, the results table and the planning documents of the repository file them are listed in Supplementary Table~18. The evaluation protocol and all splits were committed to the repository before the first main-table run. Pipeline checks had scored a reference U-Net and the released M2MRF checkpoints on the IDRiD and DDR test sets before that; the field-of-view rule and the lesion-size boundaries were then revised once, from training images only. After the freeze, the training amount and the selection rule were set from training loss and validation curves, and no test metric was used. One exception is disclosed in Supplementary Note~2: the source recipe of the reference U-Net was revised after the test score of a smoke run had been seen; this affects the source-recipe track only.

\subsection{Data collection}
\paragraph{Training sources.} IDRiD\cite{porwal2018idrid} provides 81 images of $4288\times2848$ pixels with pixel masks for MA, HE, EX and SE; the official split is 54 training and 27 test images. A missing mask file means that the image has no lesion of that class. DDR\cite{li2019ddr} provides 757 images with the official split 383\,/\,149\,/\,225.

\paragraph{Independent test sources.} TJDR\cite{mao2023tjdr} contains 561 images from one hospital: 257 standard-field images ($2048\times2048$, 35--50$^\circ$) and 304 ultra-widefield images ($3912\times3912$, 133$^\circ$), annotated with polygons by two ophthalmologists independently and adjudicated by a third; the four classes are mutually exclusive. The two fields are treated as two datasets. MAPLES-DR\cite{maplesdr2024scientificdata} adds pixel labels for 198 MESSIDOR\cite{decenciere2014messidor} images (45$^\circ$; $1440\times960$ to $2304\times1536$ pixels): seven retinologists drew MA, HE, EX and SE at $1500\times1500$ on the field-of-view crop, MA, HE and EX starting from a network's pre-segmentation and SE from scratch; the labels are placed back on the original image, and only the official test split of 60 images is used. The FGADR Seg-set\cite{zhou2021fgadr} contains 1,842 images of $1280\times1280$ pixels with masks for the four classes and, as extension labels, intraretinal microvascular abnormalities (159 images) and neovascularization (49). Retinal-Lesions\cite{wei2020retinallesionslesionnet} (release 20191227) contains 1,593 images from EyePACS, cropped to the field of view and resized to $896\times896$, from 1,505 patients, annotated with ellipses and polygons by at least three of 45 graders. DiaRetDB1 v2.1\cite{kauppi2007diaretdb1} contains 89 images of $1500\times1152$ pixels (50$^\circ$) in which four experts independently marked regions (4,135 circles, 763 ellipses, 1,036 polygons), each with a confidence level.

\paragraph{Granularity tiers.} Sources are assigned to four tiers: fine (IDRiD, DDR, TJDR standard field, TJDR ultra-widefield, MAPLES-DR), mixed (FGADR), coarse (Retinal-Lesions) and region consensus (DiaRetDB1). Results are averaged only within a tier.

\paragraph{Label mapping.} Only the four classes are mapped; other classes are stored as extension labels and do not enter the four-class mean. TJDR: palette 1, 2, 3, 4 to EX, HE, MA, SE. FGADR: HardExudate, Hemohedge, Microaneurysms, SoftExudate. Retinal-Lesions: hard\_exudate, retinal\_hemorrhage, microaneurysm, cotton\_wool\_spots; vitreous and preretinal hemorrhage are not merged into HE. DiaRetDB1: Hard\_exudates, Haemorrhages, Red\_small\_dots, Soft\_exudates. MAPLES-DR: Microaneurysms, Hemorrhages, Exudates, CottonWoolSpots; its uncertain red lesions (MA or HE) and uncertain bright lesions (EX, SE or drusen) are ignored for the two classes they could belong to, and drusen and neovascularization are extension labels.

\subsection{Model collection}
The literature search was frozen at 30 September 2026. A fixed query (three concepts in title or abstract: diabetic retinopathy; lesion, microaneurysm, exudate, hemorrhage or cotton wool; segmentation) was run on PubMed, Europe PMC and OpenAlex for 2015 onwards. It identified 1,389 records, of which 948 passed the Boolean gate and 764 remained after de-duplication by title; 112 entered the full-text queue. Screening of a priority list (27 DOIs, one dataset lead and one further paper found when a mis-assigned DOI was corrected) required colour fundus photographs, pixel-level segmentation of at least two of the four classes and an obtainable public dataset; it included 24 methods and 2 datasets and excluded 3 records. Together with a snowballed corpus the ledger holds 86 records. Only the authors' own repositories count as official code.

\rn{n.models} model variants were registered in three tiers (\rn{n.tier1}, \rn{n.tier2} and \rn{n.tier3}). Each has a reproducibility level: R0, no code; R1, model definition only; R2, inference code and weights; R3, training code; R4, training code, weights, splits and environment. Each is assigned to evaluation tracks: A, author weights; B, unified retraining of the official code; C, re-implementation from the paper, marked in the tables; a re-implementation fixes every unstated choice in a written list of assumptions before any experiment and is not tuned to the published numbers, and its parameter count is compared with the published one and reported. Methods whose repository holds the model but no training loop (WSRFNet) or whose code cannot run on current hardware (PMCNet, TensorFlow 1.13) are retrained under Track B with the loss and optimizer the paper states, written into the wrapper and disclosed as such. One method with official code and weights, HRDecoder, was missed by the frozen search and added on 14 October 2026 after a targeted check for code-available methods; the addition is recorded as a protocol deviation. Authors of methods without code received a fixed e-mail on 3 October 2026 asking for code, weights and splits, with a deadline of four weeks; where code exists without weights, the weights are requested separately. Non-responses and failures remain in the denominator of the reproducibility funnel.

\subsection{Preprocessing}
\paragraph{Unified layout.} Every source is exported to the same layout: the image at original resolution with the bytes of the release, one binary mask per class, a field-of-view mask and the split lists. A missing mask file is an all-negative image for that class; the image stays in the split.

\paragraph{Field of view.} The field-of-view (FOV) mask is the largest connected component of pixels with luminance $0.299R+0.587G+0.114B>10$, with holes filled. An Otsu threshold was abandoned after it removed the darker half of the retina in several training images. With the fixed threshold, the largest share of lesion pixels lying more than 2 pixels outside the mask on the training images is $2.3\times10^{-5}$ (IDRiD hemorrhages) and 0 on DDR. For the independent sources one rule is added: if the four $40\times40$ corner blocks are uniform (1st to 99th percentile within 5 grey levels), the threshold is the larger of 10 and the corner median plus 3. It changes one of 4,085 images, a DiaRetDB1 image with a uniform grey background.

\paragraph{Canvas.} For training and inference the image is cropped to the FOV and resized so that its long side is 1,440 pixels. Predictions are mapped back to the original resolution by bilinear interpolation, set to zero outside the FOV and stored as 16-bit images (probability = value\,/\,65535). All metrics are computed offline from these maps.

\paragraph{Source-specific handling.} FGADR masks have anti-aliased edges with values 1--254; a pixel is positive at $\geq128$, and the number of pixels dropped is recorded per class. In Retinal-Lesions, value 255 is lesion and value 127, which the authors use for annotations that do not follow the grading guideline, is written to a per-class ignore region that is skipped by all metrics and, when this source is used for training, by the loss. In DiaRetDB1, circles, ellipses and polygons are rasterized at pixel centres; for one expert and class the highest confidence per pixel is kept, with High = 1.0, Medium = 0.75 and Low = 0.25; the four experts are averaged and a pixel is positive at $\geq0.75$, the threshold of the original protocol. With this rule six of the eight positive-image counts of the original paper are matched exactly and two differ by one image (the preparation report in the repository); the union of all expert marks is kept as a lenient ground truth for sensitivity analysis.

\paragraph{Exclusive training targets.} Models with a softmax head need one label per pixel. Overlapping pixels are written in the order EX, HE, SE, MA, later classes overwriting earlier ones. This removes less than 0.5\% of the positive pixels of any class in the training sets (at most 0.34\% at original resolution). Evaluation always uses the four original, possibly overlapping masks.

\subsection{Construction of training, validation and test sets}
Official splits are never changed. New splits use the generator PCG64 with seed 20261003, and the lists were committed before any training on the source concerned. Where a split is stratified, quotas per stratum follow the largest-remainder method. IDRiD: 10 of the 54 official training images, drawn at random from the sorted identifiers, are held out for validation (44\,/\,10\,/\,27); author-recipe retraining uses all 54. DDR: official 383\,/\,149\,/\,225. TJDR: the official test set is kept and 15\% of the official training images of each field are drawn for validation, giving 172\,/\,30\,/\,55 (standard) and 209\,/\,37\,/\,58 (ultra-widefield). FGADR: stratified by file-name suffix and DR grade, 30\% test and 15\% of the remainder for validation (1,096\,/\,193\,/\,553). Retinal-Lesions: split by patient, stratified by the higher grade of the two eyes (951\,/\,164\,/\,478). DiaRetDB1: official 28\,/\,61; the 28 training images are used only to choose thresholds.

\subsection{Test scenarios and leakage audit}
The internal test is the official test set of the training source. Independent tests come from sources excluded as a whole from training and from internal testing; conclusions rest on them, as in the template study, where independent databases were held out and similar sequences were removed\cite{lu2026tcrepitopenaturemethods}. Every main-table checkpoint is applied to the test splits of the five independent sources. Ultra-widefield images are predicted at their native resolution with sliding windows, because matching pixels per degree to the 1,440-pixel canvas would require upsampling; a version resized to 1,440 pixels is reported for comparison. The per-degree scale is defined by the field of view of the training source; for a run whose training sources do not fix one (the factor cells trained on the standard-field TJDR images, on FGADR, or on merged sources), only the 1,440-pixel version is run and stands for the ultra-widefield result. For author weights, a source that the original model was trained on is removed from its independent tests. Near-duplicate images are searched with perceptual hashes and feature similarity between sources, between the splits of a source, and between the pretraining sets of foundation models and the test sets; an independent test image that duplicates a training image is removed and listed (Extended Data Fig.~3).

\subsection{Evaluation of the original models}
\paragraph{Released checkpoints (author weights).} Released checkpoints are run in a current framework; for M2MRF, whose repository requires a framework version that does not support the cards used here, the forward pass was checked by reproducing the repository numbers with the authors' function. Each checkpoint is reported with three numbers: the paper, the authors' evaluation function, and this study's protocol. The authors' function for M2MRF evaluates the whole image, overwrites overlapping labels in the order EX, HE, SE, MA and computes AUPR from 11 thresholds.

\paragraph{Author-recipe reruns.} M2MRF and HRDecoder are retrained with the authors' pipelines, augmentation and training set (all 54 IDRiD training images; three M2MRF seeds and one HRDecoder run, plus a first M2MRF run with the wrong effective batch that is reported separately) and scored at the last iteration without checkpoint selection, with the authors' function and with this study's protocol. The registered plan had a DDR run and HACDR-Net's recipe as well; these were cut when the scope was fixed, and the second and third M2MRF seeds were run when cards were free. For M2MRF: binary Dice loss (smoothing $10^{-5}$); SGD with learning rate 0.01, momentum 0.9 and weight decay $5\times10^{-4}$; polynomial decay with power 0.9; 40,000 iterations on IDRiD and 60,000 on DDR; random crops of $960\times1440$ after resizing; photometric distortion; HRNet-W48 initialized from ImageNet. For HRDecoder: the authors' segmentor and loss, images resized to $2880\times1920$ with scale jitter 0.5--2, random $1920\times1920$ crops, flips and rotations of $\pm45^\circ$, SGD 0.01 with 1,500 warm-up steps and linear decay over 20,000 steps, four images per step; their four-process synchronized batch normalization becomes one process with four images. The authors' runner evaluates on the test split every 2,000 iterations and keeps the best checkpoint; the rerun reports the last step, and what that test-split selection would have returned is given separately.

\subsection{Unified retraining}
\paragraph{What is shared.} Data, splits, the 1,440-pixel canvas, augmentation, the training amount in pixels, checkpoint selection, inference geometry and evaluation are the same for every model. Each model keeps its loss, optimizer, learning rate, weight decay, schedule shape and effective batch from its source recipe (Table~\ref{tab:recipes}); the schedule is stretched over the shared number of steps.

\paragraph{Training amount.} One equivalent epoch draws as many patch pixels as the canvases of the training split contain, so that a model trained on $224\times224$ patches draws about 18 times more patches per epoch than one trained on $960\times960$ patches:
\begin{equation}
\text{steps per epoch}=\left\lceil \frac{1}{B}\left\lceil \frac{\sum_i H_i W_i}{h\,w} \right\rceil \right\rceil ,
\end{equation}
where $H_i\times W_i$ is the canvas of training image $i$, $h\times w$ the patch and $B$ the effective batch. The number of equivalent epochs was fixed per dataset by a budget probe: U-Net and M2MRF were trained once (seed 0) at twice a provisional amount, and only training loss and validation \mAUPR{} were read. The validation curve was smoothed with a centred three-point mean; if the smoothed maximum of any probe of a dataset lay in the last 10\% of the steps, the amount was doubled. This gave 1,000 equivalent epochs for IDRiD and 400 for DDR. The rule was revised twice after the validation curves had been seen and never after a test metric had been read; all versions of the rule and their outcomes are given in Supplementary Note~2.

\paragraph{Patches and inference.} Models with a fixed input size (Swin-Unet and RETFound with UPerNet, 224; FCT, 384; CLC-Net, TC-Net and SAMed, 512; SAM-Adapter and MedSAM, 1,024) are trained on patches of that size cut from the canvas and predict with sliding windows of the same size at an overlap of 0.5; HRDecoder keeps its authors' windowed inference at 960. All other models are trained on $960\times960$ patches and predict on the whole canvas (Table~\ref{tab:recipes}).

\paragraph{Augmentation.} Random scaling by 0.5--2 before cropping, horizontal flip with probability 0.5, rotation by multiples of 90$^\circ$, and photometric distortion with probability 0.5 (brightness $\pm32$, contrast and saturation 0.5--1.5, hue $\pm18^\circ$). The randomness of a sample depends only on the seed and the sample index.

\paragraph{Checkpoint selection.} A run is validated 20 times at equal intervals and the checkpoint with the highest validation \mAUPR{} inside the FOV is scored; the last step is kept for a sensitivity analysis. A model whose best checkpoint falls in the last two validations is flagged as possibly under-trained. The test split is not read during training.

\paragraph{Reproducibility.} A resume file stores model, optimizer, step, selection state, random-number state and data order; a run resumed after an interruption is bit-identical to an uninterrupted one. Each run writes its environment (host, card, driver, framework, seed, commit at process start, parameter count, wall time) and its recipe with every deviation.

\paragraph{Rare-class collapse.} The reference U-Net uses a binary cross-entropy plus a sigmoid Dice term computed per image and class with smoothing 1. On a patch without positives this term is minimized by an all-zero output with a gradient of about 1 per pixel, whereas on a patch with a few positive pixels the gradient per pixel is about 2\,/\,(number of positive pixels + 1). Under patch training 83\% of DDR patches and 67\% of IDRiD patches contain no soft exudate. In the budget probe the soft exudate channel collapsed to zero on IDRiD, and on DDR the microaneurysm, hemorrhage and soft exudate channels did. With a smoothing of $10^{-5}$ the gradient of this term on an empty patch is small while the channel's output is far from zero, and a Dice pooled over the batch (nnU-Net) is not affected. The U-Net's Dice term was therefore pooled over the batch per class with smoothing $10^{-5}$, as in nnU-Net; the revision was decided from the training and validation curves of a diagnostic run (validation \mAUPR{} 0.569 on IDRiD and 0.544 on DDR against 0.127 before) after the registered recipe's test numbers had been read, and both versions are reported, the registered one in Supplementary Table~6 (Supplementary Note~2). Before any foundation model was trained the same rule was registered for SAM2-UNet, whose weighted IoU term is computed per image and class with smoothing 1, and for MedSAM, whose Dice term is computed per image and class with squared predictions and smoothing $10^{-5}$: if the validation AUPR of a rare class stays at its prior level throughout training, the term is pooled over the batch and the model is retrained from its initial weights. The rule applied to SAM2-UNet on DDR (microaneurysm and soft-exudate validation AUPR 0.02--0.04 and 0.05--0.10 from the first to the last evaluation) and to MedSAM on DDR (microaneurysm 0.000--0.004 and soft exudate 0.002 at every evaluation, their prior levels, while hemorrhage and exudate trained); both models were retrained with the term pooled over the batch (SAM2-UNet to 0.664 and 0.360 \mAUPR{} on IDRiD and DDR, MedSAM to 0.642 and 0.413), the main tables show the retrained runs, and the collapsed runs are kept as the registered recipe in the supplementary tables.

\paragraph{nnU-Net.} nnU-Net is run through its own entry point in 2D at original resolution, with the same splits written into its split file and exclusive labels composed in the same overwrite order; softmax probabilities are exported, mapped to the original image and scored with the same evaluation. The default configuration and the residual-encoder preset M are both run. Its preprocessing is its own and is marked as such.

\paragraph{Source-recipe track.} A second track, each model trained with its source recipe after minimal adaptation to the shared data, canvas and evaluation (seed 0), was registered to measure what moving a recipe to this task gives; it is not an architecture ranking, since the recipes differ about 67-fold in the number of images seen on IDRiD. Six of its runs were started and the track was stopped for compute on 2026-10-11 before any of them finished; no number from it is reported.

\subsection{Foundation models}
Four kinds of use are reported separately. (a) Zero-shot: automatic masks of SAM and SAM~2, scored once as class-agnostic ``any lesion'' proposals and once with each mask assigned the class it overlaps most; text prompts for SAM~3 with three wordings per class. (b) Oracle prompts: boxes or points generated from ground-truth connected components, with 1, 5, 10, 20 or all components per image and class; MedSAM, which accepts boxes only, belongs here. (c) Parameter-efficient fine-tuning: SAMed, SAM-Adapter and SAM2-UNet, trained on patches cut from the canvas (512, 1,024 and 960 pixels; Table~\ref{tab:recipes}). (d) Encoder plus decoder: MedSAM fully fine-tuned, and RETFound with a UPerNet head. To isolate the effect of ophthalmic pretraining, two pairs of ViT-L encoders fixed before any was trained are run on DDR with the same UPerNet head, recipe and training amount: RETFound's masked-autoencoder encoder against the ImageNet masked autoencoder it was initialised from, and RETFound's DINOv2 encoder against DINOv2 itself (one seed each; the pair differences are the pre-specified comparisons). After the DDR result of the first pair had been read, the three control encoders were also trained on IDRiD under the same head and recipe, so that both pairs exist on both training sources; this addition is reported as such. Settings (c) and (d) follow the unified protocol.

\subsection{Evaluation of factors}
Each experiment changes one factor and keeps the unified protocol otherwise; every cell is trained at the pixel budget of DDR with the DDR recipe of its model and one seed (seed 0), and differences smaller than twice the median seed standard deviation of the main table are reported without interpretation. \emph{Training source}: the five representative models (U-Net, nnU-Net, M2MRF, SAM2-UNet and RETFound with UPerNet) are trained on the standard-field TJDR images and on FGADR and tested on all eight test sets, which with the main table gives four training sources. \emph{Annotation granularity}: U-Net and M2MRF on six training sets at an equal number of patch pixels, namely DDR, FGADR, Retinal-Lesions, DDR with FGADR, DDR with Retinal-Lesions, and DDR with 383 FGADR images drawn once (list committed before training); the DDR and FGADR rows are the main-table and training-source runs. \emph{Input resolution}: M2MRF and RETFound with UPerNet on DDR with a long side of 512 or 1,440 pixels (the main table) and with patches at native resolution; the HRDecoder case study of the main text trains that model at its authors' 2,880-pixel canvas and 1,920-pixel window with the equivalent epochs of the unified protocol counted at that canvas. \emph{Training amount}: U-Net and M2MRF on 10\%, 50\% and 100\% of the DDR training set (nested subsets drawn once) and on IDRiD, DDR and standard-field TJDR merged. \emph{Loss}: U-Net on DDR with cross-entropy, batch Dice, their sum (the main table) and the per-image Dice of the registered recipe, with the architecture fixed. \emph{Learning rate}: HRNet-W48 and M2MRF on DDR with the learning rate multiplied by 1/3 and 3, judged on validation \mAUPR{} only. An extra vessel-map input and further resolution and training-amount levels of the registered plan were cut for compute and are listed as future work.

\subsection{Metrics}
\paragraph{Pixel level.} For each class, the positive and negative pixels of all test images inside the FOV (and outside the ignore region of that class) are pooled, and AUPR is computed exactly over the 16-bit probability levels of the stored maps; this equals the average precision
\begin{equation}
\mathrm{AUPR}=\sum_{k}\left(R_k-R_{k-1}\right)P_k ,
\end{equation}
with precision $P_k$ and recall $R_k$ at the $k$-th distinct threshold. \mAUPR{} is the mean over the four classes. Dice and intersection over union are reported at a fixed threshold of 0.5 and at the per-class threshold that maximizes Dice on the validation set. For an image without positive pixels of a class the per-image Dice is missing, not zero.

\paragraph{Lesion level.} Lesions are 8-connected components. A ground-truth lesion is detected if any predicted positive pixel overlaps it; a predicted component is a true positive if it overlaps any ground-truth pixel of the class. Precision, recall and F1 are pooled over the split. Lesion size is the equivalent diameter of the ground-truth component divided by the equivalent diameter of the FOV; the boundaries between small, medium and large, 0.005 and 0.011, are the medians of microaneurysms and hemorrhages on the 44 IDRiD and 383 DDR training images, rounded up to 0.001. Validation and test images were not used to set them.

\paragraph{Annotation style.} The labelling style of each source (Fig.~\ref{fig:sources}a) is measured on the prepared masks of all its images, every image once whatever its split: per class, the 8-connected components are counted and their equivalent diameter is divided by the long side of the FOV bounding box (which for a circular field equals the FOV diameter); ignore regions are removed before counting, and DiaRetDB1 is measured on the $\geq0.75$ consensus the protocol evaluates. Reported are the mean number of components per image, over all images of the source, and the median relative diameter over all components.

\paragraph{DiaRetDB1.} The primary result follows the official image-level protocol: the image score of a class is the maximum probability inside the FOV, and the area under the ROC curve and the weighted error rates are reported. Pixel-level AUPR against the consensus regions is given separately and marked as region ground truth.

\paragraph{Authors' implementation.} The 11-threshold AUPR is reported next to the exact value for all models, only to explain distances to published numbers.

\paragraph{Case study of one model's geometry.} HRDecoder trains its backbone at half the resolution of a crop that is itself two thirds of the authors' 2,880-pixel canvas; the unified protocol halves the canvas and keeps that ratio. Three settings separate the recipe from the budget and the canvas: the authors' recipe retrained (the author-recipe rerun), the unified protocol at twice its training amount (for HRDecoder and, as a control, HRNet-W48 on DDR, trained on cards taken out of the queue), and the unified protocol on the authors' canvas with the authors' 1,920-pixel window and equivalent epochs counted at that canvas (four times the pixels per step). All three were registered before any of them was trained (docs/paper/14, note 46 of the repository).

\paragraph{Ensembles.} Probability maps of the main-table models trained on one source are averaged pixel by pixel in five compositions fixed before any ensemble was scored: all models; the three and the five best by validation \mAUPR{}; the best model of each family; and each family on its own. The member list is the 18 main-table rows scored when the compositions were registered; MedSAM, scored afterwards, is not a member. Members are chosen on validation data only; the ensemble's own validation maps set its thresholds; the test sets and the scoring are those of the single models. Ensembles are reported as a reference below the single models and do not enter the rank tests.

\paragraph{Reference model.} As a starting point for users of the benchmark, two models are trained once on the merged training splits of the three fine-grained sources (IDRiD, DDR and the standard-field TJDR set; 599 images) with the recipe, canvas and pixel budget of their DDR main-table run, validated on DDR, and their probability maps are averaged as in the ensembles. The two were chosen before any merged run by their DDR validation \mAUPR{} among the main-table models (RETFound with UPerNet, 0.643; SAMed, 0.632); the pair, the budget rule and the validation source were registered on 2026-10-23. The reference model is scored on every test set like the ensembles (thresholds from the DDR validation split of the averaged maps), is not a ranked row, and is released with its weights and maps and with a command that applies it to new images (\texttt{python -m bench.reference predict}: field-of-view detection by the rule above, the 1,440-pixel canvas, each member's sliding-window inference, the averaged probability map at the original resolution, and masks at the validation thresholds).

\paragraph{Human reference.} On DiaRetDB1 each of the four experts in turn is scored as if it were a model: the expert's own markings, with the confidence mapping above, are the probability map, and the consensus of the other three (mean confidence $\geq0.75$) is the ground truth; thresholds come from the official training list as for the models, and the four values of each metric are reported as a band next to the DiaRetDB1 column. A map with four confidence levels gives a four-point precision--recall curve, so this band is a lower bound on what an expert's pixel-level AUPR would be with a continuous score. Lesion-level markings by three ophthalmologists on DDR and MAPLES-DR test images are pre-registered and will be reported in a Supplementary Note when complete.

\paragraph{Results by DR grade.} For the test sets that carry a DR grade per image (DDR, FGADR, Retinal-Lesions, MAPLES-DR), every score is also computed within grades 0--1 and within grades 2--4 from the stored per-image histograms, with the thresholds of the full test; ungradable images enter neither stratum. Image-level detection summarizes the maximum probability inside the FOV against the presence of the class in the ground truth.

\paragraph{Label revision.} Refined IDRiD\cite{chankhachon2026refinedidrid} re-annotates the IDRiD images as one $1024\times1024$ grey-level label per image (one class per pixel: background, vitreous hemorrhage, retina, fovea, vessels, optic disc, hard exudates, IRMA, hemorrhages, neovascularisation, cotton-wool spots and microaneurysms) on a crop of the retinal region. For each of the 27 test images the crop is located in the original image by a coordinate-descent search over the four edges of the box, started from our FOV box, under two resampling hypotheses (the box stretched to the square, or padded to a square and resized uniformly), and accepted only when the normalised cross-correlation between the resampled original and the release reaches 0.97 (all 27: padded, 0.9988--0.9996). The official masks and the FOV are resampled through the recovered geometry with nearest-neighbour interpolation; cotton-wool spots are read as soft exudates and the other refined classes are ignored. Each IDRiD-trained checkpoint's original-resolution test maps are moved through the same geometry and scored by the protocol evaluator against both label sets, with the thresholds of its own validation run; no model is retrained or re-run. The per-class differences between the two label sets (pixels, components, intersection over union) are reported with the scores.

\subsection{Statistics}
Confidence intervals are 95\% percentile intervals from 2,000 bootstrap replicates over test images (seed 0); each replicate pools the histograms of the drawn images and recomputes AUPR, and a replicate without positive pixels is dropped; the interval of \mAUPR{} averages the four class values of the same replicates and drops a replicate in which any class has no positive pixel. Point estimates use the histograms at full 16-bit resolution and bootstrap replicates use per-image histograms coarsened to 4,096 bins. Two models are compared by a paired bootstrap on AUPR with shared image indices, and by a two-sided Wilcoxon signed-rank test on per-image Dice with Holm correction. Seed mean and s.d.\ are reported separately from bootstrap intervals. Rankings are compared across test sets with the Friedman test and the Nemenyi critical difference\cite{demsar2006statistical},
\begin{equation}
\mathrm{CD}=q_{\alpha}\sqrt{\frac{k(k+1)}{6N}} ,
\end{equation}
for $k$ models and $N=8$ test sets (two internal and six independent). Scores are the seed means of Table~\ref{tab:main}; in an independent block the score of a model is the mean over its IDRiD-trained and DDR-trained checkpoints; the DiaRetDB1 block uses the official image-level AUC on DiaRetDB1 (mean over the four classes). A version restricted to the fine-grained test sets is given in the repository's analysis notebook; the distribution of ranks under the image bootstrap\cite{maierhein2018rankings} is left for the final version. Agreement between published and retrained rankings is measured by Kendall's $\tau$.

\subsection{Time and resource consumption}
Training ran on NVIDIA A100 cards with 40\,GB in a node of eight, one task per card, with a queue that assigns tasks from longest to shortest expected duration. Evaluation of author weights, the first author-recipe rerun, data export and all metric computation ran on a workstation (RTX 4060 Ti, 16\,GB). Where a card cannot hold the effective batch of a recipe, gradients are accumulated and the deviation is recorded. Parameter count, floating-point operations, peak training memory, GPU hours and inference time per image are reported for every model (Extended Data Fig.~10).

\subsection{Data availability}
IDRiD, DDR and TJDR are public. FGADR and Retinal-Lesions are available from their authors on request. The FGADR licence forbids redistribution of the images and of any derived data; Retinal-Lesions is distributed under application terms, and we handle it in the same way. DiaRetDB1 is public. The MAPLES-DR labels are public (CC BY 4.0); the MESSIDOR images they annotate are available from ADCIS on request. We release all split lists, the protocol files and, for IDRiD, DDR, TJDR, MAPLES-DR and DiaRetDB1, the probability maps of every model (the maps hold no image data; DiaRetDB1 is distributed for research use). For FGADR and Retinal-Lesions we release split lists and code only.

\subsection{Code availability}
The benchmark code, including data export, training under the unified protocol, inference, the evaluation module and the scripts that produce every figure and table of this paper, is available at \url{https://github.com/skylynf/retiseg}. Official model repositories are referenced as pinned submodules and are not modified.

\FloatBarrier
\begin{table*}[t]
\centering\tablefont
\caption{\textbf{What each model keeps under the unified protocol.} Data, splits, the 1,440-pixel canvas, augmentation, the training amount in pixels, checkpoint selection and evaluation are shared (Methods). Loss, optimizer, learning rate, weight decay, schedule shape and effective batch are taken from each model's source recipe; the schedule is stretched over the shared number of steps. Models with a fixed input size train on patches of that size and predict with half-overlapping windows; all others train on $960\times960$ patches and predict on the whole canvas. nnU-Net runs through its own entry point and keeps its own plan.}
\label{tab:recipes}
\begin{adjustbox}{max width=\textwidth}
\begin{threeparttable}
\begin{tabular}{@{}l l l l l c c l c c l r@{}}
\toprule
\colh{Model} & \colh{Initialization} & \colh{Head} & \colh{Loss} & \colh{Optim.} & \colh{LR} & \colh{WD} & \colh{Schedule} & \colh{Batch} & \colh{Patch} & \colh{Inference} & \colh{Steps} \\
\midrule
\famchip{general}\; U-Net & none & 4 sigmoid & BCE + Dice$^a$ & Adam & 1e-4 & 0 & constant & 4 & 960 & whole canvas & 21,000 \\
\famchip{general}\; DeepLabv3 & ImageNet, ResNet-101 & 4 sigmoid & BCE & SGD & 7e-3 & 1e-4 & poly 0.9 & 16 & 960 & whole canvas & 6,000 \\
\famchip{general}\; HRNet-W48 & ImageNet & 4 sigmoid & BCE & SGD & 1e-2 & 5e-4 & poly & 12 & 960 & whole canvas & 7,000 \\
\famchip{general}\; Swin-Unet & ImageNet, Swin-T & 5 softmax & 0.4 CE + 0.6 Dice & SGD & 1e-2 & 1e-4 & poly 0.9 & 24 & 224 & sliding window & 63,000 \\
\famchip{general}\; FCT & none & 4 sigmoid & BCE & Adam & 1e-3 & 0 & constant & 4 & 384 & sliding window & 129,000 \\
\famchip{auto}\; nnU-Net & none & 5 softmax & CE + batch Dice & SGD & 1e-2 & 3e-5 & poly 0.9 & 2 & 1024$\times$1536$^c$ & sliding window & 250,000$^c$ \\
\famchip{auto}\; nnU-Net ResEnc-M & none & 5 softmax & CE + batch Dice & SGD & 1e-2 & 3e-5 & poly 0.9 & 13$^c$ & 384$\times$640$^c$ & sliding window & 250,000$^c$ \\
\famchip{fundus}\; M2MRF & ImageNet, HRNet-W48 & 4 sigmoid & binary Dice & SGD & 1e-2 & 5e-4 & poly 0.9 & 4 & 960 & whole canvas & 21,000 \\
\famchip{fundus}\; HACDR-Net & none & 5 softmax & weighted CE$^b$ & AdamW & 6e-5 & 1e-2 & poly 1.0, warm-up & 1 & 960 & whole canvas & 83,000 \\
\famchip{fundus}\; HRDecoder & ImageNet, HRNet-W48 & 4 sigmoid & binary Dice + 0.1 Dice (HR crops) & SGD & 1e-2 & 5e-4 & poly 0.9 & 4 & 960 & sliding window & 21,000 \\
\famchip{fundus}\; PMCNet & ImageNet, EfficientNet-B0 & 5 softmax & CE & Adam & 1e-4 & 0 & poly 0.9 & 1 & 960 & whole canvas & 83,000 \\
\famchip{fundus}\; WSRFNet & ImageNet, EfficientNet-B1 & 4 sigmoid & binary Dice, per loop$^d$ & SGD & 1e-2 & 5e-4 & poly 0.9 & 4 & 960 & whole canvas & 21,000 \\
\famchip{fundus}\; CLC-Net$^{\rm C}$ & none & 5 softmax & Dice + BCE (cls.\ heads)$^e$ & Adam & 3e-4 & 0 & step (1/3, 5/6) & 16 & 512 & sliding window & 73,000 \\
\famchip{fundus}\; TC-Net$^{\rm C}$ & none & 5 softmax & CE & Adam & 5e-4 & 1e-5 & constant & 4 & 512 & sliding window & 73,000 \\
\famchip{foundation}\; SAMed & SAM ViT-B, LoRA r=4 & 5 softmax & 0.2 CE + 0.8 Dice & AdamW & 5e-3 & 0.1 & warm-up, poly & 24 & 512 & sliding window & 13,000 \\
\famchip{foundation}\; SAM-Adapter & SAM ViT-H, adapters & 4 sigmoid & BCE + IoU & AdamW & 2e-4 & 1e-2 & stepped cosine & 2 & 1,024 & sliding window & 37,000 \\
\famchip{foundation}\; SAM2-UNet & SAM 2 Hiera-L, adapters & 4 sigmoid & BCE + weighted IoU$^f$ & AdamW & 1e-3 & 5e-4 & stepped cosine & 12 & 960 & whole canvas & 7,000 \\
\famchip{foundation}\; MedSAM (full FT) & MedSAM ViT-B, all weights & 4 sigmoid & Dice + BCE & AdamW & 1e-4 & 1e-2 & constant & 2 & 1,024 & sliding window & 37,000 \\
\famchip{foundation}\; RETFound + UPerNet & RETFound ViT-L (MAE) & 4 sigmoid & BCE + 0.4 aux. & AdamW & 4.7e-4 & 0.05 & warm-up, cosine & 24 & 224 & sliding window & 63,000 \\
\bottomrule
\end{tabular}
\begin{tablenotes}[flushleft]\scriptsize
\item Batch is the effective batch (gradient accumulation where one card cannot hold it). Steps is for IDRiD at 1,000 equivalent epochs; one equivalent epoch draws as many patch pixels as the training canvases contain, so steps = $\lceil\lceil$canvas pixels\,/\,patch pixels$\rceil$\,/\,batch$\rceil$ per epoch. DDR uses 400 equivalent epochs (for example 83,600 steps at batch 4 and 960-pixel patches; Supplementary Table~3 lists every run). $^{\rm C}$, re-implemented from the paper. Every run's recipe record, with its deviations from the source recipe, is in the repository.
\item $^a$Per-image, per-class sigmoid Dice with smoothing 1 in the registered recipe; the main-table runs pool the Dice term over the batch (Methods). The other Dice losses use $10^{-5}$. $^b$IDRiD: class-weighted cross-entropy; DDR: the same plus the authors' NALoss with weight 0.2. $^c$nnU-Net plans patch and batch per dataset (IDRiD plans shown; DDR: default $1024\times1536$ at batch 2, ResEnc-M $512\times512$ at batch 12) and trains 1,000 epochs of 250 steps on both sources; its learning rate, weight decay and polynomial schedule are its defaults. $^d$The paper's loss and optimizer, written into the wrapper because the repository has no training loop. $^e$Local and contextual Dice plus binary cross-entropy of the two classification heads, as described in the paper. $^f$Registered per-image term; the main-table run pools it over the batch (Methods).
\end{tablenotes}
\end{threeparttable}
\end{adjustbox}
\end{table*}



{\small
\bibliographystyle{unsrtnat}
\bibliography{refs}
}

\clearpage
\section{Extended Data}
\renewcommand{\figurename}{Extended Data Fig.}
\setcounter{figure}{0}

\noindent Nine Extended Data figures are planned (the comparison with the source recipes moved to Supplementary Table~4 when that column was cut back to the runs already finished). Three of them need no images and are drawn below from the same data table as the main figures; the others are listed with their content and data source.

\begin{center}\small\sffamily
\begin{tabular}{@{}l l l l@{}}
\toprule
\textbf{Extended Data Fig.} & \textbf{Content} & \textbf{Source} & \textbf{Status} \\
\midrule
1 & Example images and labels of each source; annotation style & prepared datasets; label audit & needs images \\
2 & Field-of-view masks and canvas scale per source & FOV reports & needs images \\
3 & Leakage audit: near-duplicates between sources and splits & hash and feature audit & audit not run \\
4 & Sensitivity to the metric implementation (thresholds, FOV, overlap) & evaluation switches & with Fig.~\ref{fig:literature}d,e \\
5 & Per-image distributions and threshold sensitivity & per-image histograms & after E2 \\
6 & (withdrawn; source recipes are Supplementary Table~4) & E2b & -- \\
7 & Learning-rate sweep (DDR) & E2c & drawn \\
8 & Loss with the architecture fixed (U-Net) & E9 & drawn \\
9 & DiaRetDB1 official protocol; ultra-widefield at two scales & E3a & after E3a \\
10 & Computational cost & E10 & drawn \\
\bottomrule
\end{tabular}
\end{center}

\setcounter{figure}{6}
\begin{figure}[h!]
\centering
\includegraphics[width=\textwidth]{\figfile{ed07_lr_sweep}}
\caption{\textbf{Learning-rate sweep.} Validation \mAUPR{} on DDR of one general (HRNet) and one fundus-specific (M2MRF) network with the learning rate of the recipe multiplied by 1/3, 1 and 3. Kendall $\tau$ compares the order of the models at the recipe value with their order when each uses its best of the three values.}
\end{figure}

\setcounter{figure}{7}
\begin{figure}[h!]
\centering
\includegraphics[width=\textwidth]{\figfile{ed08_loss}}
\caption{\textbf{Loss with the architecture fixed.} \textbf{a}, \mAUPR{} of U-Net on the DDR test set with cross-entropy, batch Dice, cross-entropy with batch Dice (the main-table recipe) and cross-entropy with a per-image Dice term (the registered recipe, which collapses on rare classes); seed 0. \textbf{b}, Per-class change in AUPR against cross-entropy.}
\end{figure}

\setcounter{figure}{9}
\begin{figure}[h!]
\centering
\includegraphics[width=\textwidth]{\figfile{ed10_cost}}
\caption{\textbf{Computational cost.} \textbf{a}, \mAUPR{} on the IDRiD test set against the training time of one seed on IDRiD (40\,GB A100, one task per card); marker area, number of parameters. \textbf{b}, Parameters, trainable parameters under the unified protocol, floating-point operations for one 1,440-pixel canvas, peak training memory and inference time per image.}
\end{figure}

\clearpage
\section{Supplementary information}
\noindent\textbf{Supplementary tables.}
\begin{enumerate}[leftmargin=*,label=\textbf{\arabic*}.,itemsep=1pt]
\item Datasets: source, camera, field of view, resolution, annotation procedure, splits, licence; DiaRetDB1 positive-image counts against the original paper.
\item Models: all \rn{n.models} registered variants with tier, code, weights, reproducibility level, tracks and environment.
\item Reported numbers extracted from the literature (one row per paper, baseline, dataset, class and metric), with split, resolution, metric implementation and provenance.
\item Recipe and deviations of every run; source recipes (B1) against the unified protocol for the runs that completed.
\item Per-class results of every model on every test set: AUPR with bootstrap interval, 11-threshold AUPR, Dice and intersection over union at both thresholds, lesion-level precision, recall and F1 by lesion size.
\item Split lists.
\item Computational cost.
\item Refined IDRiD: the crop geometry recovered for each of the 27 test images, the label differences per class, and every IDRiD-trained model's per-class AUPR against the refined and the official labels through the same geometry.
\end{enumerate}
\noindent\textbf{Supplementary notes.} One note per model (source, reproducibility level, environment, changes, differences from the paper); the budget probe with every version of the rule; protocol details.


\end{document}
